\documentclass[11pt,a4paper]{article}
\usepackage[margin=23mm]{geometry}
\usepackage[T1]{fontenc}
\usepackage{lmodern,microtype,amsmath,amssymb,amsthm,mathtools,enumitem,xcolor,booktabs}
\usepackage[colorlinks=true,linkcolor=blue!45!black,citecolor=blue!45!black,urlcolor=blue!45!black]{hyperref}
\newtheorem{theorem}{Theorem}[section]
\newtheorem{lemma}[theorem]{Lemma}
\newtheorem{proposition}[theorem]{Proposition}
\newtheorem{corollary}[theorem]{Corollary}
\theoremstyle{remark}\newtheorem{remark}[theorem]{Remark}
\setlength{\parskip}{3pt}
\setlist{nosep}
\allowdisplaybreaks[1]
\title{Thresholds, spectral geometry, and persistent radiation\\in mass-critical dispersive equations}
\author{Shuxing Yang}
\date{8 October 2026}
\begin{document}
\maketitle
\begin{abstract}
We develop five connected results for critical focusing dispersive equations.
For homogeneous real-even elliptic symbols of degree $m>1$, a sequence of
negative-energy traveling waves minimizing mass is compact after scaling
and translation, except for a finite splitting into zero-energy exact
profiles; no diffuse remainder survives. Negative energy forces at least
three negative directions in the full real action Hessian, and momentum
response gives stronger branch-dependent bounds. A strictly convex
density-norm argument unifies phase-kernel simplicity for local powers
and positive-type Hartree attraction and yields a uniform nonlocal
traveling-profile mass gap. For same-order dispersions interpolating
between two homogeneous limits, we determine both variational thresholds,
classify the rational model's upper endpoint, and prove a sharp
lower-threshold selection law by a variational limit on the optimizer
manifold. Finally, any bounded global zero-momentum negative-energy BNLS
orbit must sustain exterior gradient at distances proportional to time.
These are profile, variational, spectral-form, and necessary dynamical
results; they do not constitute a general classification of near-critical
time-dependent solutions.
\end{abstract}
Conditional stability-index statements retain their spectral assumptions; the quadratic-form conclusions require no such assumptions.

\section*{What has advanced}
The original physical question concerns pure BNLS,
$iu_t=\Delta^2u-|u|^{8/N}u$, and the relationship between negative
energy and collapse without radial symmetry. The present continuation
studies the structures that distinguish persistent coherent motion from
collapse, and the changes in those structures under modified dispersion
or nonlocal attraction.

\begin{enumerate}
\item \textbf{The pure traveling-wave threshold.}
When negative-energy exact profiles exist, define their least possible
mass by $M_{\mathrm{tr}}$. A normalized minimizing sequence either
converges to a negative-energy minimizer or splits strongly into finitely
many zero-energy profiles with common parameters. If
$M_{\mathrm{tr}}<2M_c$, the latter alternative has only one component.
The theorem does not decide the endpoint sign or the value of
$M_{\mathrm{tr}}$.
\item \textbf{A spectral requirement at every mass.}
Negative energy forces a three-dimensional negative plane in the full
action Hessian. On a smooth isotropic branch with momentum parallel to
velocity it forces at least $N+2$ negative directions. The distinction
between near-ground and excited branches is thus expressed by precise
spectral conditions, rather than by an assumption that every coherent
wave has the same sign of energy.
\item \textbf{A nonlocal extension with a common proof mechanism.}
The simplicity of the phase kernel follows from equality in a strictly
convex density norm and sharp critical interpolation. For positive-type
Hartree interactions this gives the complete analytic input for a
uniform negative-energy profile mass gap, even with arbitrary admissible
subhomogeneous dispersion additions.
\item \textbf{A complete variational model and selected scales.}
For the rational same-order interpolation, the constrained energy has
two exact mass thresholds and three different upper-endpoint behaviors.
Near the lower threshold, the energy, kinetic scale, and frequency have
explicit leading constants. The next low-frequency moment selects among
all principal GN optimizers; uniqueness is unnecessary.
\item \textbf{A stronger necessary dynamical alternative.}
A forward-global uniformly $H^2$-bounded pure BNLS solution with
$E<0$ and $P=0$ must maintain nonzero time-averaged exterior gradient
at a linear spatial scale. Gradient localization, including an averaged
version, is incompatible with those assumptions. This does not exclude
persistent splitting or radiation.
\end{enumerate}

The common idea is to use the geometry of the density attraction to
control phase directions, the response to dilation to identify energy
thresholds, and the velocity--momentum response to distinguish coherent
branches. The actual evolution introduces a separate transport question;
a static profile theorem cannot supply that missing time-dependent step.

\section{Analytic foundations and conventions}
\label{sec:r10-foundations}
The Fourier transform is unitary, $D=-i\nabla$, and
$P(u)=\operatorname{Im}\int\overline u\nabla u=\int\xi|\widehat u|^2$.
The real Hilbert pairing is the real part of the complex $L^2$ pairing.
Throughout the local-power arguments, $N\ge1$, $m>1$, $s=m/2$,
$r=m/N$, and $p=2+2r$. A principal symbol $a$ is real, even,
$m$-homogeneous, $C^{N+2}$ off zero, and satisfies
$c|\xi|^m\le a(\xi)\le C|\xi|^m$.
Set $A(u)=\int a|\widehat u|^2$, $M(u)=\|u\|_2^2$, and
$E(u)=A(u)/2-\|u\|_p^p/p$. The sharp critical mass $M_c$ is defined by
\begin{equation}\label{eq:r10-basic-GN}
 \|u\|_p^p\le\frac p2\left(\frac{M(u)}{M_c}\right)^r A(u).
\end{equation}
The traveling-wave convention $u(t,x)=e^{i\omega t}U(x-vt)$ gives
$[a(D)+\omega+iv\cdot\nabla]U-|U|^{2r}U=0$.

\begin{lemma}[Normalized optimizer compactness]
The inequality \eqref{eq:r10-basic-GN} is attained. If
$M(f_n)\to M_c$, $A(f_n)\to1$, and $\|f_n\|_p^p\to p/2$, then
some translated subsequence converges strongly in $H^s$ to an optimizer.
The normalized optimizer set is compact modulo translations and phases.
\end{lemma}
\begin{proof}
Sharp Sobolev interpolation gives a finite positive best constant, and
amplitude and spatial scaling normalize a maximizing sequence as stated.
Bounded energy-space sequences that vanish locally in mass converge to
zero in $L^p$: use a bounded-overlap cover by unit balls, local Sobolev
interpolation, and $2<p<2N/(N-m)$ when $N>m$ (all finite $p$ otherwise).
Thus translations give a nonzero weak limit $Q$ and almost-everywhere
convergence. Write $M_1=M(Q)>0$, $A_1=A(Q)>0$, $M_2=M_c-M_1$,
and $A_2=1-A_1$. Quadratic splitting and Brezis--Lieb give
\[
 M_c^r\le M_1^r A_1+M_2^r A_2,\qquad A_1+A_2=1.
\]
If both masses are positive, the right side is strictly below $M_c^r$.
Consequently $M_1=M_c$ and then $A_1=1$. Weak convergence plus
convergence of both quadratic norms gives strong $H^s$ convergence.
The same proof applies to a sequence of optimizers. Constant phase is
proved by the density-norm lemma below; this compactness argument itself
does not need a positivity-preserving kinetic operator.
\end{proof}

\begin{lemma}[Weak-profile regularity and Pohozaev identities]
Every weak $H^s$ local-power profile belongs to $L^\infty\cap H^{m+1}$
and satisfies
\begin{equation}\label{eq:r10-basic-pohozaev}
 v\cdot P=2mE,\qquad \omega M=rA+(2m-p)E.
\end{equation}
No spatial moment is required.
\end{lemma}
\begin{proof}
Choose $\Lambda$ so large that
$a(\xi)-v\cdot\xi+\Lambda\asymp1+|\xi|^m$.
The shifted resolvent maps $L^q$ to $H^{m,q}$, $1<q<\infty$,
by the Mihlin multiplier theorem. For $N>m$, the starting exponent
is $q_0=2N/(N-m)$. Splitting the nonlinearity into $|U|>1$ and
$|U|\le1$, resolvent Sobolev improvement iterates the reciprocal
exponent as $\beta_{j+1}=(1+2r)\beta_j-r$ with
$\beta_0=(1-r)/2$. The ideal iterates
$\beta_j=1/2-(r/2)(1+2r)^j$ cross the boundedness threshold after
finitely many steps. Inputs have exponent greater than one; intermediate
exponents may be weakened to avoid endpoints. The linear contribution
has the stronger improvement $\beta\mapsto\beta-r$ and is handled
separately. When $N=m$ start at any sufficiently large finite exponent;
when $N<m$ the initial space already embeds in $L^\infty$.
Now the nonlinearity is in $L^2$, giving $U\in H^m\subset H^1$.
The $C^1$ map $z\mapsto|z|^{2r}z$, with bounded derivative on the
range of $U$, sends $U$ into $H^1$; one more resolvent application
gives $H^{m+1}$.

Here are the cutoff details for the dilation calculation. Let
$f_R(x)=x\chi(x/R)$ with $\chi$ compactly supported and equal to
one near zero, and let $L_R=f_R\cdot\nabla+\operatorname{div}f_R/2$.
The Fourier kernel of $[a(D),L_R]$ is, up to the transform constant,
\[
 \frac i2(a(\xi)-a(\eta))(\xi+\eta)\cdot\widehat f_R(\xi-\eta).
\]
Splitting comparable and noncomparable frequencies gives
\[
 |a(\xi)-a(\eta)||\xi+\eta|
 \le C|\xi-\eta|(|\xi|+|\eta|)^m.
\]
After division by $\langle\xi\rangle^s\langle\eta\rangle^s$,
the convolution kernel is bounded by
$C|k|\langle k\rangle^s R^{N+1}|\widehat f(Rk)|$, with a
uniform $L^1$ norm for $R\ge1$. Hence the commutators are uniformly
bounded $H^s\to H^{-s}$. Schwartz approximation and the homogeneous
symbol identity give
$\langle U,[a(D),L_R]U\rangle\to mA(U)$.
The drift commutator tends to $-v\cdot P$; its order is one and
$s>1/2$. In the nonlinear pairing integration by parts and dominated
convergence give
\[
 \operatorname{Re}\int |U|^{p-2}\overline U L_RU
 \longrightarrow N(1/2-1/p)\|U\|_p^p=\frac m p\|U\|_p^p.
\]
Testing the profile with $L_RU$ therefore gives $mE=v\cdot P/2$.
Testing with $U$ gives $A+\omega M-v\cdot P-\|U\|_p^p=0$;
combining these two equations proves \eqref{eq:r10-basic-pohozaev}.
\end{proof}

The same regularity proof holds for an individually fixed nonnegative
order-$m$ addition $b$ with the stated Mihlin derivative bounds.
Constants need not be uniform in that addition or in $v$.
For nonhomogeneous minimizers, differentiation of their scalar
mass-preserving dilation energy is legitimate directly by dominated
convergence and minimality; differentiability of the dilation orbit
as an $H^s$-valued map is not being assumed.

\section{Compactification of the negative-energy traveling-wave threshold}
\label{sec:r10-threshold}

Fix $N\ge1$, $m>1$, $s=m/2$, $r=m/N$, and $p=2+2r$. Let
$a$ be a real, even, positive elliptic $m$-homogeneous symbol, of class
$C^{N+2}$ away from zero. Throughout this section,
\[
 M(U)=\|U\|_2^2,\quad A(U)=\int a(\xi)|\widehat U(\xi)|^2\,d\xi,
 \quad E(U)=\tfrac12A(U)-\tfrac1p\|U\|_p^p,
 \quad P(U)=\int\xi|\widehat U(\xi)|^2\,d\xi.
\]
Let $M_c$ be the sharp mass in
\begin{equation}\label{eq:r10threshold-GN}
 \|f\|_p^p\le \frac p2\left(\frac{M(f)}{M_c}\right)^r A(f).
\end{equation}
An exact traveling profile is a nonzero $U\in H^s$ satisfying
\begin{equation}\label{eq:r10threshold-profile}
 [a(D)+\omega+iv\cdot\nabla]U-|U|^{p-2}U=0.
\end{equation}
The cutoff Pohozaev identities established for this equation are
\begin{equation}\label{eq:r10threshold-pohoz}
 v\cdot P(U)=2mE(U),\qquad
 \omega M(U)=rA(U)+(2m-p)E(U).
\end{equation}
Assume the set of its negative-energy traveling profiles is nonempty,
and define
\[
 M_{\mathrm{tr}}=\inf\{M(U):U\text{ satisfies }
 \eqref{eq:r10threshold-profile}\text{ for some }(\omega,v),\ E(U)<0\}.
\]
Then $M_c\le M_{\mathrm{tr}}<\infty$ by
\eqref{eq:r10threshold-GN}. The earlier phase argument strengthens this
to $M_{\mathrm{tr}}>M_c$, but the compactification proof below does not
need that improvement, optimizer uniqueness, or real nondegeneracy.

\begin{lemma}[Group velocity without a spatial moment]
\label{lem:r10threshold-velocity}
Every profile \eqref{eq:r10threshold-profile} satisfies
\begin{equation}\label{eq:r10threshold-velocity}
 vM(U)=\int\nabla a(\xi)|\widehat U(\xi)|^2\,d\xi.
\end{equation}
In particular, with
$C_a=\sup_{\xi\ne0}|\nabla a(\xi)|/a(\xi)^{(m-1)/m}$,
\begin{equation}\label{eq:r10threshold-vbound}
 |v|\le C_a\left(\frac{A(U)}{M(U)}\right)^{(m-1)/m}.
\end{equation}
For $a(\xi)=|\xi|^4$, $C_a=4$.
\end{lemma}
\begin{proof}
The profile is a critical point of
$S=E+\omega M/2-v\cdot P/2$. Choose real
$\chi\in C_c^\infty$ equal to $x_j$ near zero and put
$\chi_R(x)=R\chi(x/R)$. Differentiating
$S(e^{i\theta\chi_R}U)$ at $\theta=0$ gives
\[
 2\operatorname{Re}\langle a(D)U,i\chi_RU\rangle
 =\int v\cdot\nabla\chi_R\,|U|^2.
\]
The right side tends to $v_jM(U)$. The commutator form on the left
has Fourier kernel proportional to
\[
 (a(\xi)-a(\eta))R^{N+1}\widehat\chi(R(\xi-\eta)).
\]
The mean-value bound
$|a(\xi)-a(\eta)|\le C|\xi-\eta|(|\xi|+|\eta|)^{m-1}$,
and a division into comparable and noncomparable frequency sizes,
show that after division by
$\langle\xi\rangle^s\langle\eta\rangle^s$ the kernel is bounded by
\[
 C|\xi-\eta|\langle\xi-\eta\rangle^s
 R^{N+1}|\widehat\chi(R(\xi-\eta))|.
\]
Its convolution $L^1$ norm is uniform for $R\ge1$ by the substitution
$q=R(\xi-\eta)$. Thus the commutator forms are uniformly bounded on
$H^s$. For Schwartz $U$, passage to the limit is justified by
$\chi_RU\to x_jU$ in $L^2$, and Fourier differentiation gives
$\int\partial_ja\,|\widehat U|^2$. Density and the uniform form bound
extend this to every $H^s$ profile. H\"older's inequality then gives
\[
 |v|M(U)\le C_a A(U)^{(m-1)/m}M(U)^{1/m}.
\]
The proof never assumes differentiability of
$\eta\mapsto e^{i\eta\cdot x}U$ in $H^s$.
\end{proof}

\begin{theorem}[Finite zero-energy compactification]
\label{thm:r10threshold-compact}
Let $U_n$ be negative-energy exact traveling profiles with
$M(U_n)\to M_{\mathrm{tr}}$. Apply mass-preserving scaling so that
$A(U_n)=1$, and scale their frequencies and velocities accordingly.
After passage to a subsequence, precisely one of the following
alternatives holds.
\begin{enumerate}
\item There are translations $x_n$ and an exact traveling profile
$U_*$ such that $U_n(\cdot+x_n)\to U_*$ strongly in $H^s$, with
\[
 M(U_*)=M_{\mathrm{tr}},\qquad A(U_*)=1,\qquad E(U_*)<0.
\]
The profile parameters converge to those of $U_*$.
\item $E(U_n)\to0$. There are an integer
$1\le J\le\lfloor M_{\mathrm{tr}}/M_c\rfloor$, nonzero profiles
$V_1,\ldots,V_J$, and pairwise separating translations $x_n^j$ such that
\begin{equation}\label{eq:r10threshold-decomp}
 \left\|U_n-\sum_{j=1}^J V_j(\cdot-x_n^j)\right\|_{H^s}\longrightarrow0.
\end{equation}
All profiles solve \eqref{eq:r10threshold-profile} with the same
limiting parameters $(\omega_*,v_*)$, and
\begin{gather}
 E(V_j)=0,\qquad M(V_j)\ge M_c,\qquad
 \sum_jM(V_j)=M_{\mathrm{tr}},\qquad \sum_jA(V_j)=1,
 \label{eq:r10threshold-sums}\\
 \omega_*=\frac{r}{M_{\mathrm{tr}}}>0,\qquad
 v_*\cdot P(V_j)=0,\qquad \omega_*M(V_j)=rA(V_j).
 \label{eq:r10threshold-endpointparameters}
\end{gather}
\end{enumerate}
In particular, no diffuse mass or kinetic remainder survives at this
threshold. The second alternative is a statement about a sequence of
exact profiles, not a time-dependent soliton-resolution theorem.
\end{theorem}
\begin{proof}
The normalization is
$U_\lambda(x)=\lambda^{N/2}U(\lambda x)$ with
$\lambda=A(U)^{-1/m}$. The parameters become
$\omega_\lambda=\lambda^m\omega$ and
$v_\lambda=\lambda^{m-1}v$, and $E(U_\lambda)=\lambda^mE(U)$.
By \eqref{eq:r10threshold-GN},
\[
 -\frac12\left[\left(\frac{M(U_n)}{M_c}\right)^r-1\right]
 \le E(U_n)<0.
\]
Lemma~\ref{lem:r10threshold-velocity} bounds $v_n$, and
\eqref{eq:r10threshold-pohoz} bounds $\omega_n$. Take a subsequence
with $E(U_n)\to e\le0$, $v_n\to v_*$, and $\omega_n\to\omega_*$.
Then
\begin{equation}\label{eq:r10threshold-omegalimit}
 \omega_*M_{\mathrm{tr}}=r+(2m-p)e.
\end{equation}

Apply the standard translation profile decomposition in $H^s$.
The exponent $p$ lies strictly between $2$ and the Sobolev endpoint.
The resulting nonzero profiles $V_j$ have separating translation
parameters, quadratic norm decoupling, and a remainder vanishing in
$L^p$ in the iterated limit:
$\lim_{J\to\infty}\limsup_{n\to\infty}\|R_n^J\|_p=0$.
No small energy-space remainder is assumed at this stage. Therefore
\begin{equation}\label{eq:r10threshold-profilebounds}
 \sum_jM(V_j)\le M_{\mathrm{tr}},\qquad
 \sum_jA(V_j)\le1,\qquad
 \lim_n\|U_n\|_p^p=\sum_j\|V_j\|_p^p.
\end{equation}
Local compactness into $L^p$ and convergence of the parameters allow
passage to the equation for each translated weak limit. Hence each
$V_j$ is an exact profile for $(\omega_*,v_*)$. This step uses no
positivity assumption on $a(\xi)+\omega_*-v_*\cdot\xi$.

If some $E(V_j)<0$, the definition of $M_{\mathrm{tr}}$ and mass
decoupling force $M(V_j)=M_{\mathrm{tr}}$ and strong $L^2$
convergence of $U_n(\cdot+x_n^j)$ to $V_j$. The differences $w_n$
are bounded in $H^s$, so \eqref{eq:r10threshold-GN} gives
$\|w_n\|_p\to0$. Interpolation, using $m>1$, gives
$\|w_n\|_{H^{1/2}}\to0$, hence $P(U_n)\to P(V_j)$.
The first identity in \eqref{eq:r10threshold-pohoz} now yields
$E(U_n)\to E(V_j)$. Convergence of the nonlinear norms then gives
$A(U_n)\to A(V_j)$. Weak convergence and convergence of the
Hilbert norms prove strong $H^s$ convergence, establishing the
first alternative.

If $e<0$, a negative-energy profile must occur. Indeed, from
\eqref{eq:r10threshold-profilebounds},
\begin{equation}\label{eq:r10threshold-energysplit}
 e=\sum_j E(V_j)+\frac12\left(1-\sum_j A(V_j)\right).
\end{equation}
The series converges because the quadratic and nonlinear norm series
converge. If every $E(V_j)\ge0$, its right side is nonnegative.

It remains to consider $e=0$ with no negative-energy profile.
All terms on the right of \eqref{eq:r10threshold-energysplit} are
nonnegative. Therefore every nonzero profile has zero energy and
$\sum_jA(V_j)=1$. Sharp GN gives $M(V_j)\ge M_c$, so there are
only finitely many such profiles, with the asserted bound on $J$.
Their kinetic energies sum to one, so $J\ge1$.
Equation~\eqref{eq:r10threshold-omegalimit} gives
$\omega_*=r/M_{\mathrm{tr}}>0$. Apply
\eqref{eq:r10threshold-pohoz} to each zero-energy profile and sum:
\[
 \omega_*\sum_jM(V_j)=r\sum_jA(V_j)=r
 =\omega_*M_{\mathrm{tr}}.
\]
Thus their masses also sum to $M_{\mathrm{tr}}$. Both quadratic
remainders vanish, proving \eqref{eq:r10threshold-decomp}. The
remaining identities are the individual Pohozaev identities.
\end{proof}

\begin{corollary}[Attainment criteria]
If a threshold-minimizing sequence satisfies
$-E(U_n)/A(U_n)\ge\eta>0$, then $M_{\mathrm{tr}}$ is attained by
a negative-energy traveling profile. If $M_{\mathrm{tr}}<2M_c$,
every normalized minimizing sequence has a strongly convergent
translated subsequence to a single profile of mass
$M_{\mathrm{tr}}$ and energy at most zero. The latter assertion
does not decide whether the endpoint energy is strictly negative.
\end{corollary}

\begin{proposition}[Mass response along a traveling branch]
\label{prop:r10threshold-branch}
Suppose $z\mapsto U_z$ is a $C^2$ branch of exact profiles at fixed
frequency $\omega>0$ and velocity $ze$, with $|e|=1$ and real
$U_0=Q$. Put $\Pi(z)=e\cdot P(U_z)$ and $\kappa=\Pi'(0)$.
Then
\begin{equation}\label{eq:r10threshold-branchidentity}
 \omega\frac{d}{dz}M(U_z)
 =\frac{(m-1)z\Pi'(z)-\Pi(z)}{m},
\end{equation}
and
\[
 E(U_z)=\frac{\kappa z^2}{2m}+o(z^2),\qquad
 M(U_z)=M(Q)+\frac{(m-2)\kappa z^2}{2m\omega}+o(z^2).
\]
If $m>2$ and $\kappa<0$, the nearby nonzero branch therefore has
both negative energy and mass below $M(Q)$; in particular
$M_{\mathrm{tr}}<M(Q)$.
\end{proposition}
\begin{proof}
For $S_z=E+\omega M/2-z e\cdot P/2$, criticality gives
$\frac{d}{dz}S_z(U_z)=-\Pi(z)/2$.
Substitute $E(U_z)=z\Pi(z)/(2m)$ from Pohozaev and rearrange to
obtain \eqref{eq:r10threshold-branchidentity}.
Since $Q$ is real, $\Pi(0)=0$. Taylor expansion and integration give
the conclusions. In the validated four-dimensional cubic BNLS
branch, $m=4$, $\omega=1$, and
$\Pi(z)=2Iz+O(z^3)$ with $I<0$, so
$M(U_z)=M(Q)+Iz^2/2+O(z^4)$ and
$M(U_z)-M(Q)=2E(U_z)+O(z^4)$.
\end{proof}

\section{Negative energy, Morse index, and the momentum response matrix}
\label{sec:r10-inertia}

This section distinguishes two statements.  The negative-index obstruction
below holds for every exact profile and requires no nondegeneracy.  Statements
about derivatives with respect to velocity apply only when a differentiable
family of profiles has been supplied.  Such a family is not presumed to exist
at every profile.

Let $a$ be a real-even positive elliptic homogeneous symbol of degree $m>1$,
let $X=H^{m/2}(\mathbb R^N;\mathbb C)$, and set
\[
 \alpha=\frac{2m}{N},\qquad
 E(U)=\frac12\int a(\xi)|\widehat U(\xi)|^2\,d\xi
       -\frac1{\alpha+2}\int|U|^{\alpha+2},\qquad
 P(U)=\operatorname{Im}\int\overline U\nabla U.
\]
Suppose that a nonzero $U\in X$ solves
\begin{equation}\label{eq:r10-profile}
 a(D)U+\omega U+i v\cdot\nabla U-|U|^\alpha U=0.
\end{equation}
The mass-critical Pohozaev identity is $v\cdot P(U)=2mE(U)$.
All pairings below are form pairings on $X$ and $X^*$; in particular
the derivative pairings are defined since $m>1$.
Our complex pairing is antilinear in its first argument, and the real
Hilbert pairing is its real part.

Define the real scalar operator and the complex-linear operator
\[
 A_U=a(D)+\omega-|U|^\alpha,\qquad
 K_U=A_U+i v\cdot\nabla.
\]
Both are selfadjoint as forms on the complex space, and $A_U$ commutes
with complex conjugation.  The Hessian of
$S_{\omega,v}=E+\omega M/2-v\cdot P/2$ is a \emph{real-linear} operator
$\mathcal L_U$, whose quadratic form is
\begin{equation}\label{eq:r10-fullhessian}
 q_{\mathcal L_U}(\eta)
 =\operatorname{Re}\langle\eta,K_U\eta\rangle
 -\alpha\int |U|^{\alpha-2}
       \bigl(\operatorname{Re}(\overline U\eta)\bigr)^2.
\end{equation}
At $U=0$ the last integrand is defined to be zero.  It is bounded by
$|U|^\alpha|\eta|^2$ and defines a continuous form even when $\alpha<2$.
The notation $n_-(\mathcal L)$ denotes the maximal dimension of a real
subspace on which its quadratic form is strictly negative.  It does not
denote the number of exponentially growing modes of the evolution.

\subsection{An unconditional negative-index obstruction}

\begin{theorem}[Two scalar directions and three action directions]
\label{thm:r10-three-directions}
If $E(U)<0$, then
\begin{equation}\label{eq:r10-index-bound}
 n_-(A_U|_{X_{\mathbb R}})\ge2,
 \qquad n_-(\mathcal L_U)\ge3.
\end{equation}
More generally these conclusions require only $v\cdot P(U)<0$ and
the profile equation; homogeneity and criticality are used solely to
infer that sign from the energy.

Consequently, a profile whose full action Hessian has negative index
at most two cannot have negative energy.  These statements are uniform
in frequency and speed and do not require proximity to a ground state.
\end{theorem}
\begin{proof}
Write $U=h+ib$ with real $h,b$.  The profile equation gives
\[
 A_Uh=v\cdot\nabla b,\qquad A_Ub=-v\cdot\nabla h.
\]
Antisymmetry of the derivative and
$P(U)=-2\langle b,\nabla h\rangle$ imply
\begin{equation}\label{eq:r10-scalar-plane}
 \langle h,A_Uh\rangle=\langle b,A_Ub\rangle
       =\frac12v\cdot P(U),\qquad
 \langle h,A_Ub\rangle=0.
\end{equation}
The functions $h,b$ are linearly independent: dependence would make
$U$ a constant phase times a real function and would force $P(U)=0$.
Thus their real span is a two-dimensional negative subspace.

For the full Hessian it is useful to retain complex notation.
Since $K_UU=0$, conjugation of the profile equation gives
\[
 K_U\overline U=2i v\cdot\nabla\overline U,
 \qquad
 \langle\overline U,K_U\overline U\rangle=2v\cdot P(U).
\]
The latter pairing is real.  Selfadjointness and $K_UU=0$ show
that all $K_U$ form pairings between $U$ and $\overline U$ vanish.
Moreover, $U$ and $\overline U$ are linearly independent over
$\mathbb C$, for the same constant-phase reason as above.

For $t\in\mathbb R$ and $z\in\mathbb C$ put
$\eta=tU+z\overline U$.  Equation \eqref{eq:r10-fullhessian} becomes
\begin{equation}\label{eq:r10-three-plane}
 q_{\mathcal L_U}(tU+z\overline U)
 =2v\cdot P(U)|z|^2
 -\alpha\int |U|^{\alpha-2}
     \left(t|U|^2+\operatorname{Re}(z\overline{U}^{2})\right)^2.
\end{equation}
If $z\ne0$ the first term is strictly negative.  If $z=0$ and
$t\ne0$, the second term equals
$-\alpha t^2\|U\|_{\alpha+2}^{\alpha+2}<0$.
The real span of $U,\overline U,i\overline U$ is therefore a
three-dimensional negative subspace.  On the full four-dimensional
real space $\operatorname{span}_{\mathbb C}\{U,\overline U\}$,
the restricted form has precisely three negative directions and the
one-dimensional kernel spanned by $iU$.
\end{proof}

\begin{remark}[What the obstruction does and does not say]
The proof is an explicit finite-dimensional compression of the
Hessian, not a computation of its full spectrum.  If the essential
spectrum reaches negative values, the variational index may be
infinite; three discrete negative eigenvalues are not being asserted.
In particular it
does not by itself prove spectral or nonlinear instability.  The same
argument applies to a real dispersion operator and a gauge-invariant
nonlinearity $f(|u|^2)u$ whenever $f'\ge0$ on the range of $|U|^2$
and $\int f'(|U|^2)|U|^4>0$: the last term in
\eqref{eq:r10-fullhessian} is then replaced by
$-2\int f'(|U|^2)(\operatorname{Re}(\overline U\eta))^2$.
The implication from negative energy to $v\cdot P<0$ must be checked
separately for such a model.  No such monotonicity is asserted for
every nonlinear correction in the preceding dilation cone.
\end{remark}

\begin{remark}[Positive-type Hartree interaction]
The full index bound also extends to the critical Hartree setting
with a positive interaction form $D$, under the Hartree form and
regularity assumptions used in this manuscript.  There the frozen
complex-linear operator still obeys $K_UU=0$ and
$K_U\overline U=2iv\cdot\nabla\overline U$, while
\[
 q_{\mathcal L_U}(\eta)=\langle\eta,K_U\eta\rangle_{\mathbb R}
 -2D\bigl(\operatorname{Re}(\overline U\eta),
          \operatorname{Re}(\overline U\eta)\bigr).
\]
Positivity of $D$ makes the correction nonpositive, and its value
on the amplitude direction is strictly negative when
$D(|U|^2,|U|^2)>0$.  The proof on
$\operatorname{span}_{\mathbb R}\{U,\overline U,i\overline U\}$
is unchanged.  A nonnegative subhomogeneous dispersion addition
preserves the implication from $E<0$ because its dilation identity
is $v\cdot P=2mE-D_b<0$.  Thus $n_-(\mathcal L_U)\ge3$
also holds in that class.  The conclusion uses positivity of the
interaction form, not merely pointwise nonnegativity of its kernel.
\end{remark}

\begin{corollary}[A continuation barrier from a ground state]
\label{cor:r10-continuation}
A real sharp mass-critical ground state has full action Hessian
negative index one.  A continuous Fredholm family of profiles starting
at such a state cannot reach a negative-energy profile while retaining
only its phase and translation kernel and keeping all other spectral
values uniformly separated from zero.  If the essential spectral gap
stays open, reaching a negative-energy profile requires a net gain of
at least two negative directions through additional zero eigenvalues,
counted with multiplicity.
\end{corollary}
\begin{proof}
At a real optimizer $Q$, the phase block is nonnegative.  On the real
mass tangent space $\{h:\langle Q,h\rangle=0\}$ the real block is
nonnegative as well: $E$ is nonnegative on the mass sphere
$M=M_c$ and vanishes at $Q$, so its constrained second variation is
nonnegative.  On $Q$ itself the real block has quadratic form
$-\alpha\|Q\|_{\alpha+2}^{\alpha+2}<0$.  Hence the full index is one.
The index is locally constant for continuous selfadjoint Fredholm
forms after quotienting a kernel of fixed dimension, as long as the
remaining spectrum is separated from zero.  Apply
Theorem~\ref{thm:r10-three-directions} at the proposed endpoint.
\end{proof}

\subsection{Additional negative directions from momentum response}

Assume here that $v\mapsto U_v$ is a $C^1$ family in $X$ of exact
profiles at fixed $\omega$, and that differentiating the profile
equation is justified in $X^*$.  Define
\[
 B=D_vP(U_v).
\]
The matrix is symmetric: if $d(v)=S_{\omega,v}(U_v)$, then
$d_v=-P(U_v)/2$ and $B=-2d_{vv}$.

\begin{theorem}[A response-refined index bound]
\label{thm:r10-response-index}
At a negative-energy member of this family,
\begin{equation}\label{eq:r10-response-index}
 n_-(\mathcal L_{U_v})\ge
 3+n_-\bigl(B|_{P(U_v)^\perp}\bigr).
\end{equation}
For an isotropic, rotationally equivariant family with
$P(U_v)=p(c)v/c$, $c=|v|>0$, this gives
\begin{equation}\label{eq:r10-isotropic-index}
 n_-(\mathcal L_{U_v})\ge N+2.
\end{equation}
Consequently such a wave is not a local energy minimum on the
joint mass and momentum level set, provided that level set is a
$C^2$ manifold near the wave.
\end{theorem}
\begin{proof}
Let $w_\zeta=D_vU_v[\zeta]$.  Profile differentiation gives
\[
 \mathcal L_{U_v}w_\zeta=-i\zeta\cdot\nabla U_v,
 \qquad
 \langle w_\eta,\mathcal L_{U_v}w_\zeta\rangle_{\mathbb R}
       =\tfrac12\eta^TB\zeta.
\]
The second formula uses
$DP(U)[w]=2\operatorname{Re}\langle w,-i\nabla U\rangle$.
The pairings of $\mathcal Lw_\zeta$ with $\overline U$ and
$i\overline U$ vanish because $\int U\partial_jU=0$.
Its pairing with $U$ is $\zeta\cdot P(U)$, which vanishes for
$\zeta\perp P(U)$.
Thus the image of a negative subspace of $B|_{P^\perp}$ is negative
and form-orthogonal to the three-dimensional space in
\eqref{eq:r10-three-plane}.  The map is injective on that subspace,
since its quadratic form is strictly negative.  The two negative
spaces have trivial intersection, proving \eqref{eq:r10-response-index}.

In the isotropic case, with $e=v/c$,
\[
 B=p'(c)e\otimes e+\frac{p(c)}c(I-e\otimes e).
\]
Negative energy gives $cp(c)=v\cdot P<0$, so $p(c)<0$ and
$B|_{P^\perp}=(p(c)/c)I$ has $N-1$ negative directions.
This proves \eqref{eq:r10-isotropic-index}.
The joint mass and momentum tangent space has codimension at most
$N+1$, so it intersects the $(N+2)$-dimensional negative subspace
nontrivially.  At a constrained critical point the corresponding
second variation of $E$ is the restriction of $\mathcal L$.
If the joint level set is a $C^2$ manifold, this negative tangent
direction is realized by a constraint-preserving curve and excludes
a local minimum.
\end{proof}

The last level-set assumption is not a hidden spectral hypothesis.
For a sufficiently regular nonzero profile its charge differentials
are independent: a real relation between them would give
$(\tau-i\zeta\cdot\nabla)U=0$, forcing its $L^2$ Fourier transform
onto a hyperplane (or forcing $U=0$ if $\zeta=0$).
The finite-dimensional implicit-function theorem then supplies the
local charge manifold in the energy space.

\subsection{An exact reduction of the charge Hessian}

Suppose now that $(\omega,v)\mapsto U_{\omega,v}$ is a $C^1$
profile family on an open parameter set, that
$d(\omega,v)=S_{\omega,v}(U_{\omega,v})$ is $C^2$, and that its
charges obey the critical scaling covariance
\begin{align*}
 M(U_{\lambda^m\omega,\lambda^{m-1}v})&=M(U_{\omega,v}),\\
 P(U_{\lambda^m\omega,\lambda^{m-1}v})&=\lambda P(U_{\omega,v})
\end{align*}
whenever the parameters are in the domain.  This follows, in
particular, if the family is closed under the usual critical dilation,
up to phase and translation.  Let $H=d''(\omega,v)$ and
$B=D_vP(U_{\omega,v})$.

\begin{theorem}[Bordered momentum formula]
\label{thm:r10-bordered}
If $\omega\ne0$, the symmetric matrices $H$ and
\begin{equation}\label{eq:r10-bordered}
 -\frac12
 \begin{pmatrix}
 (m-1)v\cdot P&P^T\\
 P&B
 \end{pmatrix}
\end{equation}
are congruent.  If $B$ is invertible, define
\begin{equation}\label{eq:r10-kappa}
 \kappa=P^TB^{-1}P-(m-1)v\cdot P.
\end{equation}
Then, writing $p(H)$ for the number of positive eigenvalues,
\begin{equation}\label{eq:r10-hessian-inertia}
 p(H)=n_-(B)+\mathbf1_{\{\kappa>0\}},
 \qquad \dim\ker H=\mathbf1_{\{\kappa=0\}}.
\end{equation}
The Schur complement of the $v$ block in $H$ is exactly
$\kappa/(2m^2\omega^2)$.
\end{theorem}
\begin{proof}
Stationarity gives $d_\omega=M/2$ and $d_v=-P/2$.
Put $D=(m\omega,(m-1)v)^T$.  Differentiating the two scaling laws
gives the exact column identity
\[
 HD=\begin{pmatrix}0\\-P/2\end{pmatrix},
 \qquad H_{vv}=-B/2.
\]
The invertible matrix
\[
 R=\begin{pmatrix}m\omega&0\\(m-1)v&I\end{pmatrix}
\]
therefore satisfies $R^THR=$ \eqref{eq:r10-bordered}.
Sylvester's law of inertia and the Schur complement of $B$ prove
\eqref{eq:r10-hessian-inertia}.
For the final assertion, the mixed block is
\[
 H_{v\omega}=-\frac{P-(m-1)Bv}{2m\omega},
\]
and the upper-left entry follows from $HD=(0,-P/2)^T$.
Substitution and cancellation give
$H_{\omega\omega}-H_{\omega v}H_{vv}^{-1}H_{v\omega}
=\kappa/(2m^2\omega^2)$.
\end{proof}

For an isotropic family $P=p(c)e$ with $p(c)p'(c)\ne0$,
\begin{equation}\label{eq:r10-elasticity}
 \kappa=\frac{p(c)^2}{p'(c)}-(m-1)cp(c).
\end{equation}
Thus charge-Hessian degeneracy occurs precisely at
$cp'(c)/p(c)=1/(m-1)$, under these nonvanishing assumptions.
This is a degeneracy criterion, not an assertion that every such
point changes dynamical stability.  The momentum and its derivative
at a fixed frequency suffice; a separate numerical differentiation
in frequency is unnecessary.

\subsection{The small-speed transition at dispersion order two}

Assume a $C^1$ branch at $\omega=1$ starts at a real standing
profile $Q$.  Assume $L_-=a(D)+1-|Q|^\alpha$ has kernel
$\mathbb RQ$ and a bounded inverse on its orthogonal complement.
Assume also that the real-part derivative of the branch at zero
vanishes after the chosen translation gauge.  One sufficient
condition for this last assertion is
$\ker L_+=\operatorname{span}\{\partial_jQ\}$ together with
the translation-orthogonal slice; neither this nondegeneracy nor
branch existence is asserted for all ground states.
Define
\[
 T_{ij}=\langle\partial_iQ,L_-^{-1}\partial_jQ\rangle.
\]
If $T$ has $k$ negative eigenvalues, then
\begin{equation}\label{eq:r10-standing-index}
 n_-(L_-)\ge k,\qquad
 n_-\bigl(\operatorname{diag}(L_+,L_-)\bigr)\ge2k+1.
\end{equation}
Indeed, on the inverse-response vectors $L_-^{-1}(\zeta\cdot\nabla Q)$
the $L_-$ form is exactly $\zeta^TT\zeta$.  Let $V$ be the resulting
$k$-dimensional negative subspace.  The form of $L_-$ is negative
on $V$ and vanishes on $Q$, with no cross term.  Since
$L_+=L_--\alpha|Q|^\alpha$, it is strictly negative on
$V\oplus\mathbb RQ$: for a nonzero $V$ component the first form is
already negative, and on a nonzero multiple of $Q$ the subtracted
potential term is strictly positive.  The two diagonal blocks then
give \eqref{eq:r10-standing-index}.
Finite strictly negative subspaces persist under form-continuous
perturbations, so this lower bound persists on sufficiently nearby
traveling profiles.  In particular an isotropic negative response
$T=I_0I_N$, $I_0<0$, requires at least $2N+1$ negative action
directions on the standing state and its nearby traveling branch.

Differentiating the imaginary equation gives
$D_{v_j}U|_{v=0}=-iL_-^{-1}\partial_jQ$, and hence
\begin{equation}\label{eq:r10-small-speed}
 P(U_{1,v})=2Tv+o(|v|),\qquad B(1,v)=2T+o(1).
\end{equation}
For a nonsingular $T$ and directions on which $v^TTv$ stays away
from zero relative to $|v|^2$, the preceding exact formulas yield
\begin{align}
 E(U_{1,v})&=\frac1m v^TTv+o(|v|^2),\label{eq:r10-energy-expansion}\\
 M(U_{1,v})&=M(Q)+\frac{m-2}{m}v^TTv+o(|v|^2),
 \label{eq:r10-mass-expansion}\\
 \kappa(1,v)&=2(2-m)v^TTv+o(|v|^2).
 \label{eq:r10-kappa-expansion}
\end{align}
For \eqref{eq:r10-mass-expansion}, integrate $d_v=-P/2$ to get
$d(1,v)=M(Q)/2-v^TTv/2+o(|v|^2)$ and use weighted homogeneity,
$d=d_\omega+(m-1)v\cdot d_v/m$ at $\omega=1$.
Only the stated $C^1$ branch and $C^2$ action regularity are used;
no fourth derivative of a noninteger power nonlinearity is assumed.

At a sharp ground state $T$ is positive definite by phase
coercivity.  For sufficiently small nonzero velocity the charge
Hessian is consequently negative definite if $m>2$, whereas it
has exactly one positive eigenvalue if $1<m<2$.
At $m=2$ this leading-order test vanishes.  General homogeneous
symbols of degree two need not be quadratic forms, so no exact
Galilean degeneracy is inferred for all such symbols.

These are charge and index statements.  Under the additional
hypotheses of a Hamiltonian index theorem---a closed Hamiltonian
linearization, finite negative index, a coercive complementary
subspace, exactly the phase and translation kernel, independent
charges, and an invertible charge Hessian---the standard index
formula reads
\begin{equation}\label{eq:r10-hamiltonian-index}
 k_r+2k_c+2k_i^-
 =n_-(\mathcal L)-p(H)
 =n_-(\mathcal L)-n_-(B)-\mathbf1_{\{\kappa>0\}}.
\end{equation}
Here the terms on the left count real unstable eigenvalues,
complex unstable pairs, and the negative index on nonzero
imaginary modes, respectively, with the conventions of the chosen
index theorem.  This last formula invokes the established
Hamiltonian index theory of Lin and Zeng~\cite{LZ}; it is not a replacement
for checking its operator hypotheses.  For example, if a
nondegenerate small-speed ground-state branch retains
$n_-(\mathcal L)=1$ and all those hypotheses hold, the $m>2$
case has index one and therefore one positive real eigenvalue.
No unconditional nonlinear stability or instability conclusion
is made for arbitrary branches.

\section{A density-norm principle and nonlocal critical attraction}
\label{sec:r10-hartree}

The phase-kernel argument has a first-order formulation that applies
to both local powers and nonlocal Hartree interactions. It avoids
second derivatives of the nonlinear density at nodal sets.

\subsection{The common variational lemma}

Let $A$ be a positive quadratic form on a complex function space,
with $A(X+iY)=A(X)+A(Y)$ for real $X,Y$, and let
$M(u)=\|u\|_2^2$. Let $\mathcal B$ be a strictly convex real normed
space containing the squared densities under consideration. Fix
$\kappa>1$ and set
\[
 \mathcal F(u)=\mathfrak f(|u|^2),\qquad
 \mathfrak f(\rho)=\|\rho\|_{\mathcal B}^{\kappa}.
\]
Assume that the sharp inequality
\begin{equation}\label{r10:density-gn}
 \mathcal F(u)\le\kappa
       \left(\frac{M(u)}{M_c}\right)^{\kappa-1}A(u)
\end{equation}
has a nonzero optimizer. Assume the relevant density norm is
Gateaux differentiable at the squared density of each optimizer.
For a real optimizer $Q$ of mass $M_c$, write its derivative as
\[
 D\mathfrak f(Q^2)[h]=\kappa\langle V_Q,h\rangle.
\]
The notation allows either a multiplication potential or a bounded
linear functional on the density space; in particular
$\langle V_Q,b^2\rangle$ must be a well-defined quadratic form.

\begin{lemma}[Simple phase kernel from density convexity]
\label{lem:r10-density}
Under these assumptions every complex optimizer is a constant phase
times a real function. For a real mass-$M_c$ optimizer $Q$, the form
\[
 \mathcal L_-(b)=A(b)+(\kappa-1)\frac{A(Q)}{M_c}M(b)
                    -\langle V_Q,b^2\rangle
\]
is nonnegative on real functions and has kernel
$\operatorname{span}_{\mathbb R}\{Q\}$.
\end{lemma}
\begin{proof}
For $u=X+iY$, the density triangle inequality, sharp interpolation,
and two-term H\"older give
\begin{align*}
 \mathcal F(u)^{1/\kappa}
 &\le\mathcal F(X)^{1/\kappa}+\mathcal F(Y)^{1/\kappa}\\
 &\le \kappa^{1/\kappa}M_c^{-(\kappa-1)/\kappa}
  \left[A(X)^{1/\kappa}M(X)^{(\kappa-1)/\kappa}
       +A(Y)^{1/\kappa}M(Y)^{(\kappa-1)/\kappa}\right]\\
 &\le \kappa^{1/\kappa}M_c^{-(\kappa-1)/\kappa}
              A(u)^{1/\kappa}M(u)^{(\kappa-1)/\kappa}.
\end{align*}
All inequalities are equalities for an optimizer. Strict convexity
in the density space forces $Y^2=cX^2$ unless a component vanishes.
Apply the same assertion after multiplication by $e^{i\pi/4}$.
If both signs of $Y/X$ occurred on sets of positive measure, the
rotated squared components would fail to be proportional; when
$c=1$, they would instead have disjoint nonzero supports.
Thus one relative sign holds almost everywhere, proving constant
phase. Here, as usual for a normed space of densities, equality in
$\mathcal B$ identifies densities almost everywhere.

Put $M=M_c$, $A=A(Q)$, $B=M(b)$ and $C=A(b)$.
For real $b$, \eqref{r10:density-gn} gives
\[
 \mathfrak f(Q^2+tb^2)
 \le\kappa(1+tB/M)^{\kappa-1}(A+tC),\qquad t\ge0.
\]
Differentiation at zero proves $\mathcal L_-(b)\ge0$.
The homogeneity of $\mathfrak f$ gives
$\langle V_Q,Q^2\rangle=\mathfrak f(Q^2)=\kappa A$,
so $\mathcal L_-(Q)=0$.

If $b\ne0$ belongs to the form kernel, the first-derivative bound
is an equality. The supporting-functional bound for a norm power is
\[
 D\mathfrak f(Q^2)[b^2]
 \le\kappa\mathfrak f(Q^2)^{(\kappa-1)/\kappa}
                   \mathfrak f(b^2)^{1/\kappa}.
\]
Use $\mathfrak f(Q^2)=\kappa A$ and
\eqref{r10:density-gn} for $b$. With $x=C/A$ and $y=B/M$,
the kernel equality implies
\[
 x+(\kappa-1)y
 \le\kappa x^{1/\kappa}y^{(\kappa-1)/\kappa}
 \le x+(\kappa-1)y.
\]
The last inequality is weighted arithmetic--geometric mean.
Hence $x=y$ and equality holds in the supporting-functional bound.
Strict convexity makes the supporting hyperplane meet the unit
sphere at one point, giving $b^2=yQ^2$. Since a nonnegative
quadratic form has a linear kernel, apply this conclusion also to
$b+Q$. Subtracting the squared identities proves $b=cQ$ with one
real constant $c$.
\end{proof}

For the local critical power, take
$\mathcal B=L^{\kappa}$ and $\kappa=1+m/N$.
For a positive-type Hartree interaction, take the Hilbert space of
densities in its interaction norm and $\kappa=2$.
Lemma~\ref{lem:r10-density} itself asserts neither compactness of
near-optimizers nor a dynamical theorem. Those require additional
analytic information, supplied for the Hartree model below.

\begin{corollary}[The local-power profile gap]
For the homogeneous local-power class of Section~\ref{sec:r10-foundations},
there is $\delta>0$ such that every negative-energy traveling profile
has mass at least $M_c+\delta$. Consequently, when the negative-energy
profile set is nonempty, $M_{\mathrm{tr}}>M_c$.
\end{corollary}
\begin{proof}
For a real normalized optimizer $Q$, the preceding lemma with
$\kappa=1+r$ gives the nonnegative phase form
$L_Q=a(D)+r/M_c-|Q|^{2r}$ with kernel $\mathbb RQ$.
Multiplication by $|Q|^{2r}$ is compact $H^s\to H^{-s}$:
H\"older with two $L^p$ factors controls its tails, and local
Rellich compactness controls its restriction to a ball. Hence
$\langle h,L_Qh\rangle\ge c_Q\|h\|_{H^s}^2$ for real $h\perp Q$.
If negative-energy profiles had masses tending to $M_c$, normalize
$A(U_n)=1$. Sharp GN then gives $E(U_n)\to0$ and convergence to a
normalized optimizer modulo translations and phase. Make its limit
$Q$ real and choose phases so $Y_n=\operatorname{Im}U_n\perp Q$.
Pohozaev gives $\omega_n\to r/M_c$, and strong $H^s$ convergence
gives convergence of $|U_n|^{2r}$ in potential form norm.
The imaginary equation and $P(U_n)=-2\int Y_n\nabla\operatorname{Re}U_n$
give the exact identity
\[
 mE(U_n)=\langle Y_n,[a(D)+\omega_n-|U_n|^{2r}]Y_n\rangle
 \ge\tfrac12c_Q\|Y_n\|_{H^s}^2\ge0
\]
for large $n$, a contradiction.
\end{proof}

\subsection{The critical Hartree class}

Fix $1<m<N$, set $s=m/2$, and let $a$ be real, even, positive elliptic,
$C^{N+2}$ away from zero, and homogeneous of degree $m$:
\[
 c|\xi|^m\le a(\xi)\le C|\xi|^m.
\]
Let $K$ be real, even, $C^1$ away from zero, homogeneous of degree
$-m$, and satisfy
\begin{equation}\label{r10:kernel-bounds}
 |K(x)|+|x||\nabla K(x)|\le C_K|x|^{-m}.
\end{equation}
Assume its interaction form
\[
 \mathcal D(f,g)=\iint K(x-y)f(x)g(y)\,dx\,dy
\]
is positive definite on the real density space in use. More
explicitly, it suffices that the Fourier transform of $K$ is a
locally integrable, strictly positive function almost everywhere;
then the Fourier representation of $\mathcal D$ gives the required
inner product and identifies zero-norm densities with zero functions.
Examples include $K(x)=|x|^{-m}$ and positive finite sums of
$|B_jx|^{-m}$ for invertible real matrices $B_j$.
Pointwise positivity of $K$ is not needed for the proof.

Write
\[
 H(u)=\mathcal D(|u|^2,|u|^2),\qquad
 A(u)=\int a(\xi)|\widehat u(\xi)|^2\,d\xi.
\]
Hardy--Littlewood--Sobolev and Sobolev interpolation give the finite
sharp inequality
\begin{equation}\label{r10:hartree-gn}
 H(u)\le\frac{2}{M_c}M(u)A(u).
\end{equation}
Indeed, with $q=4N/(2N-m)$, the interaction is bounded by
$C\|u\|_q^4$, and interpolation gives
$\|u\|_q^4\le C M(u)A(u)$.

Allow in addition any real-even symbol $b\ge0$ satisfying
\begin{equation}\label{r10:hartree-bclass}
 \xi\cdot\nabla b\le mb,\qquad
 |\partial^\beta b(\xi)|\le C_{b,\beta}|\xi|^{-|\beta|}
                    (1+|\xi|^m),\quad |\beta|\le N+2,
\end{equation}
with $b\in C^{N+2}(\mathbb R^N\setminus\{0\})$.
There is no common bound on the individual constants.
Define
\[
 B_b(u)=\int b(\xi)|\widehat u(\xi)|^2\,d\xi,
 \qquad E_b(u)=\tfrac12[A(u)+B_b(u)]-\tfrac14H(u).
\]

\begin{theorem}[Nonlocal mass gap over the dispersion cone]
\label{thm:r10-hartree-gap}
There exists $\delta>0$, depending only on $N,m,a,K$, such that
for every admissible $b$, every $\omega\in\mathbb R$ and
$v\in\mathbb R^N$, a nonzero weak $H^{m/2}$ solution of
\begin{equation}\label{r10:hartree-profile}
 [a(D)+b(D)+\omega+iv\cdot\nabla]U-(K*|U|^2)U=0
\end{equation}
with $E_b(U)<0$ satisfies $M(U)\ge M_c+\delta$.
The gap is uniform over the entire addition class, frequencies,
and velocities. No radial symmetry, spatial moment, positivity or
uniqueness of a ground state, or real-sector nondegeneracy is assumed.
\end{theorem}

\subsection{Compactness and the phase form}

Normalized maximizing sequences for \eqref{r10:hartree-gn} are compact
modulo translation and constant phase. Here are the details needed
for this nonlocal statement. Normalize $M(u_n)=M_c$, $A(u_n)=1$, and
$H(u_n)\to2$. Since $H(u_n)\le C\|u_n\|_q^4$, local vanishing in
mass is impossible by the localized Sobolev inequality. A translated
subsequence has a nonzero weak $H^s$ limit $Q$, with local strong
$L^q$ convergence. Set $r_n=u_n-Q$.
The product $Qr_n$ tends to zero in $L^{q/2}$ by local convergence
and the small $L^q$ tail of $Q$.
Moreover,
\begin{equation}\label{r10:hartree-crossvanish}
 \mathcal D(|Q|^2,|r_n|^2)\longrightarrow0.
\end{equation}
To see this, approximate $Q$ in $L^q$ by a smooth compactly supported
function. Its potential is bounded and tends to zero at infinity,
so multiplication by it is compact as a form $H^s\to H^{-s}$.
The Hardy--Littlewood--Sobolev estimate controls the approximation
error, proving \eqref{r10:hartree-crossvanish}. It also proves the
nonlocal splitting
\[
 H(u_n)=H(Q)+H(r_n)+o(1).
\]
If $M_j,A_j$ are the limiting mass and kinetic contributions of the
two pieces, \eqref{r10:hartree-gn} implies
\[
 M_c\le M_1A_1+M_2A_2,
 \qquad M_1+M_2=M_c,\quad A_1+A_2=1,\quad M_1>0.
\]
Both masses cannot be positive. Thus $M_1=M_c$, $A_1=1$, proving
strong convergence and attainment. The same argument covers
sequences with these normalization limits and all sequences of
normalized optimizers.

Take $\mathcal B$ to be the completion of real densities under
$\|\rho\|_{\mathcal B}^2=\mathcal D(\rho,\rho)$.
Lemma~\ref{lem:r10-density} shows that every optimizer has constant
phase. A real optimizer with $M(Q)=M_c$, $A(Q)=1$ satisfies
\[
 a(D)Q+\omega_*Q-(K*Q^2)Q=0,\qquad \omega_*=1/M_c.
\]
The real phase form
\[
 L_Q=a(D)+\omega_*-K*Q^2
\]
is nonnegative with kernel $\operatorname{span}_{\mathbb R}\{Q\}$.
The potential is a compact form by the same approximation used in
\eqref{r10:hartree-crossvanish}. Consequently
\begin{equation}\label{r10:hartree-phasecoercivity}
 \langle h,L_Qh\rangle\ge c_Q\|h\|_{H^s}^2
 \quad\text{for real }h\perp Q.
\end{equation}
The constants can be taken uniform over normalized optimizers:
otherwise translate a convergent sequence of optimizers and use
form continuity to contradict \eqref{r10:hartree-phasecoercivity}
for its limit. Form continuity follows from HLS, since
$u\mapsto |u|^2$ is continuous $H^s\to L^{q/2}$.

\subsection{Weak-profile regularity and identities}

Every weak profile in \eqref{r10:hartree-profile} belongs to
$L^\infty\cap H^{m+1}$. For the fixed $b,v$, choose $\Lambda$ large
so that
\[
 a(\xi)+b(\xi)-v\cdot\xi+\Lambda\asymp1+|\xi|^m.
\]
The resolvent maps $L^p$ to $H^{m,p}$ for $1<p<\infty$ by the
individual Mihlin bounds. Split $K$ into its parts inside and outside
a fixed ball. The outside potential applied to $|U|^2$ is bounded.
The inside kernel lies in $L^j$ for $1\le j<N/m$, while
$|U|^2\in L^{N/(N-m)}$. Young's inequality therefore puts the
inside potential in arbitrarily large finite $L^t$ spaces.
Fix sufficiently large $t>N/m$. Multiplication by this potential,
followed by the resolvent, improves a reciprocal integrability
exponent by $m/N-1/t>0$. Starting from the energy Sobolev exponent,
finitely many iterations give $U\in L^\infty$; intermediate
exponents may be weakened to avoid endpoints. Bounded potential
and linear terms improve at least as strongly and may be combined
using interpolation with $L^2$.

Now $K*|U|^2$ is bounded, giving $U\in H^m\subset H^1$.
Also $\nabla |U|^2\in L^1\cap L^2$, and the near/far kernel split
gives $\nabla(K*|U|^2)\in L^2+L^\infty$.
Thus $(K*|U|^2)U\in H^1$, and the resolvent gives $H^{m+1}$.

Put $\rho=|U|^2$, $f_R(x)=x\chi(x/R)$ with smooth compactly
supported $\chi=1$ near zero, and
$L_R=f_R\cdot\nabla+(\operatorname{div}f_R)/2$.
Symmetrized integration by parts gives
\begin{align}
 &\operatorname{Re}\int(K*\rho)U\,\overline{L_RU}\,dx\notag\\
 &\quad=-\frac14\iint\rho(x)\rho(y)
       [f_R(x)-f_R(y)]\cdot\nabla K(x-y)\,dx\,dy
 \longrightarrow\frac m4H(U).
 \label{r10:hartree-cutoff}
\end{align}
One can first regularize or exclude the diagonal. The omitted
boundary term vanishes because $m<N$ after symmetrization.
The uniform Lipschitz bound for $f_R$ and
\eqref{r10:kernel-bounds} dominate the integrand by
$C\rho(x)\rho(y)|x-y|^{-m}$, integrable by HLS, and justify the
limit without a spatial moment.

For clarity, the kinetic commutator used here has kernel, up to
Fourier normalization,
\[
 \frac i2[c(\xi)-c(\eta)](\xi+\eta)\cdot
                    \widehat f_R(\xi-\eta),\qquad c=a+b.
\]
The symbol assumptions give
\[
 |c(\xi)-c(\eta)||\xi+\eta|
 \le C|\xi-\eta|[1+(|\xi|+|\eta|)^m].
\]
Divide by the two energy weights and use
$\widehat f_R(k)=R^{N+1}\widehat f(Rk)$ to get an integrable kernel
majorant with a bound independent of $R$. The limiting form for
Schwartz functions is $\int\xi\cdot\nabla c\,|\widehat U|^2$;
density extends it to the energy space. The degree-one drift term
has limit $-v\cdot P(U)$, where
$P(U)=\int\xi|\widehat U|^2$ is finite because $m>1$.

Abbreviate $A=A(U)$, $B=B_b(U)$, $M=M(U)$ and $E=E_b(U)$, and set
\[
 D_b=\int[mb(\xi)-\xi\cdot\nabla b(\xi)]
                    |\widehat U(\xi)|^2\,d\xi\ge0.
\]
Testing by $L_RU$, taking the limit, and testing separately by $U$
gives
\begin{equation}\label{r10:hartree-identities}
 v\cdot P=2mE-D_b,
 \qquad \omega M=A+B+(2m-4)E-D_b.
\end{equation}
Write $U=X+iY$ with real components and put
$\widetilde\omega=\omega+D_b/M$.
The imaginary profile equation, together with
$P(U)=-2\int Y\nabla X$, yields the exact shifted identity
\begin{equation}\label{r10:hartree-shift}
 \begin{aligned}
 mE={}&\langle Y,[a(D)+\widetilde\omega-K*|U|^2]Y\rangle
       +B_b(Y)\\
     &+D_b\left(\frac12-\frac{M(Y)}M\right),\qquad
 \widetilde\omega M=A+B+(2m-4)E.
 \end{aligned}
\end{equation}

\subsection{Proof of the uniform nonlocal gap}

\begin{proof}[Proof of Theorem~\ref{thm:r10-hartree-gap}]
Negative energy and \eqref{r10:hartree-gn} require $M>M_c$.
If the gap failed, take profiles with $M(U_n)\to M_c$ and negative
energies, allowing arbitrary admissible $b_n,\omega_n,v_n$.
The critical scaling
\[
 U_\lambda(x)=\lambda^{N/2}U(\lambda x),\quad
 b_\lambda(\xi)=\lambda^m b(\xi/\lambda),\quad
 \omega_\lambda=\lambda^m\omega,\quad
 v_\lambda=\lambda^{m-1}v
\]
preserves the class, mass, and profile equation, and multiplies
energy by $\lambda^m$. Normalize $A(U_n)=1$.
Writing $d_n=M(U_n)/M_c-1$, the sharp bound gives
\[
 0\le B_{b_n}(U_n)<d_n,\qquad
 -d_n/2\le E_n<0,\qquad H(U_n)\longrightarrow2.
\]
Compactness yields, after translations and phases, strong
$H^s$ convergence to a real normalized optimizer $Q$.
A further phase makes $Y_n=\operatorname{Im}U_n$ orthogonal to $Q$.
Equation~\eqref{r10:hartree-shift} gives
$\widetilde\omega_n\to1/M_c$ even if $D_{b_n}$ is unbounded.
HLS gives convergence of the Hartree potentials in form norm.
By \eqref{r10:hartree-phasecoercivity}, for large $n$,
\[
 \langle Y_n,[a(D)+\widetilde\omega_n-K*|U_n|^2]Y_n\rangle
       \ge c\|Y_n\|_{H^s}^2.
\]
The other terms in \eqref{r10:hartree-shift} are nonnegative,
since $M(Y_n)/M(U_n)\to0$. Its right side is therefore
nonnegative, contradicting $mE_n<0$.
\end{proof}

The conclusion extends a local attraction mechanism to a nonlocal
density interaction. It is a profile theorem, not a classification
of arbitrary global trajectories. The positive-density geometry is
an explicit hypothesis; arbitrary matrix couplings, non-even
symbols, and indefinite interaction kernels are not covered.

\section{Two dispersion thresholds and the selection of critical profiles}
\label{sec:r10-symbols}

The loss of a profile gap under a small nonhomogeneous perturbation has a
more precise variational description.  We prove two distinct mass
thresholds, classify the upper endpoint for a full interpolation family,
and determine the scale and selected profiles as the lower threshold is
approached.  The selection argument requires neither uniqueness of the
homogeneous optimizers nor nondegeneracy of their real linearization.
These are constrained variational statements; they do not classify
arbitrary solutions of the homogeneous evolution.

\subsection{Symbols and the two thresholds}

Fix $N\ge1$, $m>1$, $s=m/2$, $r=m/N$, and $p=2+2r$.  Let $a_0$ and
$a_\infty$ be real, even, positive elliptic symbols homogeneous of degree
$m$, continuous on the sphere and $C^{N+2}$ off zero.  Suppose
$a_\infty-a_0$ is also positive elliptic.  Use the notation
\[
 M(f)=\|f\|_2^2,\qquad A_j(f)=\int a_j(\xi)|\widehat f(\xi)|^2\,d\xi,
 \qquad E_j(f)=\tfrac12 A_j(f)-\tfrac1p\|f\|_p^p
 \quad(j=0,\infty).
\]
The sharp homogeneous masses $M_j$ are defined by
\begin{equation}\label{eq:r10-sharp-gn}
 \frac1p\|f\|_p^p\le
 \frac12\left(\frac{M(f)}{M_j}\right)^r A_j(f).
\end{equation}
In particular, $M_\infty>M_0$: if
$a_\infty\ge(1+\eta)a_0$, sharpness gives
$M_\infty\ge(1+\eta)^{1/r}M_0$.

Write the full symbol as $k=a_0+b$, with $b$ real, even, and continuous
off zero.  Assume that, for some $q>0$ and constants $c,C>0$,
\begin{equation}\label{eq:r10-bounds}
 c\min\{|\xi|^{m+q},|\xi|^m\}
 \le b(\xi)\le C\min\{|\xi|^{m+q},|\xi|^m\},
\end{equation}
and that, uniformly for $\theta\in S^{N-1}$,
\begin{equation}\label{eq:r10-high-limit}
 R^{-m}k(R\theta)\longrightarrow a_\infty(\theta)
 \qquad(R\longrightarrow\infty).
\end{equation}
Thus $k$ is comparable with $|\xi|^m$.  Define
\[
 B(f)=\int b(\xi)|\widehat f(\xi)|^2\,d\xi,\qquad
 E_k(f)=\tfrac12[A_0(f)+B(f)]-\tfrac1p\|f\|_p^p,
 \qquad I(M)=\inf_{\substack{f\in H^s\\M(f)=M}}E_k(f).
\]

\begin{theorem}[Two homogeneous limits determine two mass thresholds]
\label{thm:r10-two-thresholds}
Under \eqref{eq:r10-bounds}--\eqref{eq:r10-high-limit}, the sharp GN
mass for the full form $A_0+B$ is exactly $M_0$, and its sharp inequality
has no nonzero optimizer.  The constrained infima satisfy
\begin{enumerate}
 \item $I(M)=0$, without attainment, for $0<M\le M_0$;
 \item $-\infty<I(M)<0$ for $M_0<M<M_\infty$. Every minimizing
 sequence has, after translations and extraction, a strong $H^s$ limit
 attaining $I(M)$. A real minimizer can be chosen;
 \item $I(M)=-\infty$ for $M>M_\infty$.
\end{enumerate}
The symbols $a_0$ and $a_\infty$ need not have the same anisotropy.
\end{theorem}

\begin{proof}
Since $A_0+B\ge A_0$, the full GN constant is at most the principal
one. If $Q$ is a principal optimizer and
$Q_\lambda(x)=\lambda^{N/2}Q(\lambda x)$, then
\[
 \lambda^{-m}B(Q_\lambda)
 =\int\lambda^{-m}b(\lambda\xi)|\widehat Q(\xi)|^2\,d\xi
 \longrightarrow0\qquad(\lambda\downarrow0)
\]
by dominated convergence and \eqref{eq:r10-bounds}. Thus the sharp
constants coincide. Since $B(f)>0$ for every nonzero $f$, equality in
the full inequality is impossible. At mass $M\le M_0$, the sharp
principal inequality gives $E_k(f)>0$. Dilating any smooth mass-$M$
trial to low frequency proves that the infimum is zero.

For every $\eta>0$, \eqref{eq:r10-high-limit} gives a finite $C_\eta$
such that $k\ge(1-\eta)a_\infty-C_\eta$. Therefore, at fixed mass $M$,
\begin{equation}\label{eq:r10-coercivity}
 E_k(f)\ge\frac12\left[1-\eta-(M/M_\infty)^r\right]A_\infty(f)
             -\frac12C_\eta M.
\end{equation}
If $M<M_\infty$, choose $\eta$ to make the bracket positive. This
proves a finite lower bound and boundedness of minimizing sequences.
If also $M>M_0$, sharpness and amplitude scaling give a real smooth
mass-$M$ trial $f$ with $E_0(f)<0$. Its low-frequency dilation satisfies
\[
 E_k(f_\lambda)=\lambda^m E_0(f)+O(\lambda^{m+q})<0
\]
for sufficiently small $\lambda$, proving $I(M)<0$.

We record the compactness argument. Amplitude scaling gives
\[
 E_k(\sqrt\vartheta f)
 =\vartheta^{p/2}E_k(f)
   +\tfrac12(\vartheta-\vartheta^{p/2})[A_0(f)+B(f)]
 \qquad(\vartheta>1).
\]
Thus $I(\vartheta L)\le\vartheta I(L)$, with
$I(\vartheta L)\le\vartheta^{p/2}I(L)<\vartheta I(L)$ if $I(L)<0$,
whenever the masses are below $M_\infty$. For a mass $M$ with
$I(M)<0$, it follows that $I(M)/M<I(L)/L$ for every $0<L<M$.
Consequently
\begin{equation}\label{eq:r10-binding}
 I(M)<I(L)+I(M-L)\qquad(0<L<M).
\end{equation}
A bounded minimizing sequence cannot vanish locally in mass: the
localized GN inequality would then imply $\|f_n\|_p\to0$, contradicting
$I(M)<0$. After translations there is a nonzero weak $H^s$ limit $Q$
and almost-everywhere convergence. The quadratic forms split by weak
Hilbert convergence, and the nonlinear integral by the Brezis--Lieb
lemma. If $0<M(Q)<M$, normalize the remainder to its limiting mass;
boundedness of its kinetic and nonlinear terms makes the energy change
$o(1)$. Splitting then contradicts \eqref{eq:r10-binding}. Hence
$M(Q)=M$. For the remainder $w_n=f_n-Q$, the principal GN inequality
gives
\[
 \|w_n\|_p^p\le C M(w_n)^r A_0(w_n)\longrightarrow0.
\]
Energy splitting now forces $A_0(w_n)+B(w_n)\to0$, giving strong $H^s$
convergence and attainment.

Realness does not require a modulus inequality for a higher-order
operator. Write $f=X+iY$, put $\vartheta=M(X)/M(f)$, and, when both
components are nonzero, set $g=X/\sqrt\vartheta$,
$h=Y/\sqrt{1-\vartheta}$. These real functions have the same mass as
$f$. Evenness of $k$ splits its form over real and imaginary parts,
while convexity gives
\[
 |f|^p\le\vartheta|g|^p+(1-\vartheta)|h|^p,
 \qquad E_k(f)\ge\vartheta E_k(g)+(1-\vartheta)E_k(h).
\]
Therefore a real minimizing sequence exists and its strong limit is real.

Finally, for $M>M_\infty$ choose a smooth mass-$M$ function with
$E_\infty(f)<0$. High-frequency dilation and
\eqref{eq:r10-high-limit} give
\[
 \lambda^{-m}E_k(f_\lambda)\longrightarrow E_\infty(f)<0
 \qquad(\lambda\longrightarrow\infty),
\]
using $k\le C|\xi|^m$ for domination. Hence $I(M)=-\infty$.
\end{proof}

For each minimizer $U$ in the open window, the multiplier equation and
its scalar identity are
\begin{equation}\label{eq:r10-multiplier}
 k(D)U+\omega U-|U|^{2r}U=0,
 \qquad \omega M=r[A_0(U)+B(U)]-pI(M)>0.
\end{equation}
If additionally $k$ is $C^{N+2}$ off zero and satisfies the individual
order-$m$ derivative estimates
\[
 |\partial^\beta k(\xi)|\le C_\beta|\xi|^{m-|\beta|},
 \qquad |\beta|\le N+2,
\]
the elliptic bootstrap gives $U\in L^\infty\cap H^{m+1}$. A real
minimizer has zero momentum and its standing-wave formula gives a
global solution directly.

\subsection{The complete isotropic upper-endpoint classification}

Consider the interpolation
\begin{equation}\label{eq:r10-interpolation}
 h=a_\infty-a_0,\qquad
 k_q(\xi)=a_0(\xi)+h(\xi)\frac{|\xi|^q}{1+|\xi|^q}
         =a_\infty(\xi)-d_q(\xi),\qquad
 d_q(\xi)=\frac{h(\xi)}{1+|\xi|^q}.
\end{equation}
Write $D_q(f)=\int d_q|\widehat f|^2$. At mass $M_\infty$,
\begin{equation}\label{eq:r10-endpoint-energy}
 E_{k_q}(f)=E_\infty(f)-\tfrac12D_q(f),\qquad E_\infty(f)\ge0.
\end{equation}

\begin{theorem}[Upper-endpoint trichotomy]
\label{thm:r10-endpoint}
For \eqref{eq:r10-interpolation}, $q<m$ implies
$I(M_\infty)=-\infty$, whereas $q>m$ implies
$-\infty<I(M_\infty)<0$, attainment by a real minimizer, and
precompactness of every minimizing sequence modulo translations.
Both assertions allow anisotropic $a_0,a_\infty$.

If $a_0(\xi)=|\xi|^m$ and
$a_\infty(\xi)=(1+\varepsilon)|\xi|^m$, then at $q=m$,
\begin{equation}\label{eq:r10-borderline-endpoint}
 I(M_\infty)=-\tfrac12\varepsilon M_\infty,
 \qquad M_\infty=(1+\varepsilon)^{N/m}M_0,
\end{equation}
and the infimum is not attained.
\end{theorem}

\begin{proof}
Suppose $q<m$ and take an $a_\infty$ optimizer $Q$ of mass
$M_\infty$, so $E_\infty(Q)=0$. Then
\[
 \lambda^{-(m-q)}D_q(Q_\lambda)
 =\int\frac{h(\xi)}{\lambda^{-q}+|\xi|^q}|\widehat Q(\xi)|^2\,d\xi
 \longrightarrow
 \int h(\xi)|\xi|^{-q}|\widehat Q(\xi)|^2\,d\xi>0.
\]
The limit is finite: $0<m-q<m$, and the integrands are dominated by
$C|\xi|^{m-q}|\widehat Q|^2$. Thus
\eqref{eq:r10-endpoint-energy} tends to $-\infty$ along this sequence.

Suppose $q>m$. Then $d_q$ is bounded and tends to zero at infinity, so
\eqref{eq:r10-endpoint-energy} supplies the finite lower bound
$-\|d_q\|_\infty M_\infty/2$. Negativity follows from the
low-frequency trial since $M_\infty>M_0$. We must prove boundedness of
minimizing sequences because \eqref{eq:r10-coercivity} degenerates at
this endpoint. If instead a subsequence satisfies
$A_\infty(f_n)=t_n^m\to\infty$, write
$f_n(x)=t_n^{N/2}V_n(t_nx)$. Then
\[
 M(V_n)=M_\infty,\quad A_\infty(V_n)=1,\quad
 E_\infty(V_n)=t_n^{-m}[E_{k_q}(f_n)+D_q(f_n)/2]\longrightarrow0.
\]
Normalized sharp-GN compactness gives, after translations,
$V_n\to Q$ strongly in $H^s$, where $Q$ is an $a_\infty$ optimizer.
The bounded multiplier and strong $L^2$ convergence permit replacement
of $V_n$ by $Q$ in
\[
 D_q(f_n)=\int d_q(t_n\xi)|\widehat V_n(\xi)|^2\,d\xi.
\]
The expression for $Q$ tends to zero by dominated convergence. Thus
$D_q(f_n)\to0$ and \eqref{eq:r10-endpoint-energy} implies
$\liminf E_{k_q}(f_n)\ge0$, contradicting the negative infimum.
Minimizing sequences are therefore bounded. Amplitude scaling proves
\eqref{eq:r10-binding} also at $M=M_\infty$, since both positive
submasses lie strictly below it. The compactness and realness arguments
of Theorem~\ref{thm:r10-two-thresholds} now apply unchanged.

In the isotropic case $q=m$,
$d_m(\xi)=\varepsilon|\xi|^m/(1+|\xi|^m)<\varepsilon$.
Every admissible nonzero function consequently has
$E_{k_m}(f)>-\varepsilon M_\infty/2$. Conversely, high-frequency
dilations of an $a_\infty$ optimizer have zero homogeneous energy and
$D_m(Q_\lambda)\to\varepsilon M_\infty$ by dominated convergence.
This proves the exact infimum and its nonattainment.
\end{proof}

For completeness, the anisotropic borderline $q=m$ cannot be replaced
by the isotropic formula. Let $\mathcal G_\infty$ be the equality set
normalized by $M=M_\infty$, $A_\infty=1$. Define
\[
 J_{\rm esc}=-\frac12\sup_{Q\in\mathcal G_\infty}
       \int h(\xi/|\xi|)|\widehat Q(\xi)|^2\,d\xi.
\]
The supremum is attained by normalized GN compactness and continuity of
this bounded multiplier form. High-frequency optimizer trials give
$I(M_\infty)\le J_{\rm esc}$. If this inequality is strict,
minimizing sequences are bounded and attainment follows. Indeed an
unbounded sequence, normalized exactly as in the preceding proof,
has a limiting homogeneous optimizer and energy liminf at least
$J_{\rm esc}$. We make no assertion about attainment in the general
anisotropic equality case.

Every minimizer of \eqref{eq:r10-interpolation}, including the attained
upper endpoint $q>m$, obeys the exact identity
\begin{equation}\label{eq:r10-exact-negative-energy}
 E_{k_q}(U)=-\frac{q}{2m}
 \int h(\xi)\frac{|\xi|^q}{(1+|\xi|^q)^2}|\widehat U(\xi)|^2\,d\xi<0.
\end{equation}
To prove it, differentiate the explicit Fourier energy of the
mass-preserving dilation at its scalar minimum $\lambda=1$. The
derivative is dominated by $C|\xi|^m$. Subtract the resulting identity
divided by $m$ from the energy. This uses no spatial moment.

\subsection{Sharp behavior at the lower threshold}

Return to the general symbol and assume additionally, uniformly on
$S^{N-1}$,
\begin{equation}\label{eq:r10-low-expansion}
 R^{-(m+q)}b(R\theta)\longrightarrow c_*(\theta)
 \qquad(R\downarrow0),
\end{equation}
where $c_*$ is positive elliptic and homogeneous of degree $m+q$.
Let
\[
 \mathcal G_0=\{Q:M(Q)=M_0,\ A_0(Q)=1,\ \|Q\|_p^p=p/2\},
 \qquad C_*(f)=\int c_*(\xi)|\widehat f(\xi)|^2\,d\xi.
\]
Assume that at least one $Q\in\mathcal G_0$ has $C_*(Q)<\infty$.
This finite-moment assumption is automatic for $0<q\le m+2$ under the
stated regularity of $a_0$: its optimizer equation and the elliptic
bootstrap give $Q\in H^{m+1}$. For larger $q$ it is an explicit
hypothesis; no unlimited smoothness of a non-polynomial nonlinearity is
presumed.

Define
\begin{equation}\label{eq:r10-selection-constants}
 \begin{gathered}
 C_{\min}=\min_{Q\in\mathcal G_0}C_*(Q),\qquad
 \beta_* =\left[\frac{m}{(m+q)C_{\min}}\right]^{1/q},\qquad
 K_* =\frac{q}{2(m+q)}\beta_*^m,\\
 \mathcal G_{\rm sel}=\{Q\in\mathcal G_0:C_*(Q)=C_{\min}\}.
 \end{gathered}
\end{equation}
These definitions are well posed. The assumed finite trial gives
finiteness; positivity follows from Jensen's inequality,
\[
 \int|\xi|^{m+q}|\widehat Q|^2
 \ge M_0^{-q/m}\left(\int|\xi|^m|\widehat Q|^2\right)^{1+q/m},
\]
and $A_0(Q)=1$. A minimizing sequence in $\mathcal G_0$ is compact in
$H^s$ modulo translations, and the nonnegative $C_*$ form is lower
semicontinuous. Hence the minimum is attained. The selected set is
compact in $H^s$ modulo translations and constant phase.

\begin{theorem}[Scale law and optimizer selection]
\label{thm:r10-selection}
Let $M\downarrow M_0$ and let $U_M$ be any minimizer of $I(M)$. Set
\begin{equation}\label{eq:r10-normalization}
 \alpha=\frac{M}{M_0},\qquad d=\alpha^r-1,\qquad
 \lambda_M=\left[\frac{A_0(U_M)}{\alpha}\right]^{1/m},\qquad
 U_M(x)=\sqrt\alpha\,\lambda_M^{N/2}V_M(\lambda_Mx).
\end{equation}
Then $M(V_M)=M_0$, $A_0(V_M)=1$, and
\begin{align}
 I(M)&=-\alpha K_*d^{1+m/q}(1+o(1)),\label{eq:r10-energy-law}\\
 \lambda_M/d^{1/q}&\longrightarrow\beta_*,\label{eq:r10-scale-law}\\
 \operatorname{dist}_{H^s}(V_M,\mathcal G_{\rm sel})&\longrightarrow0,
       \label{eq:r10-shape-law}\\
 1-\frac2p\|V_M\|_p^p&=o(d),\label{eq:r10-deficit-law}\\
 \frac{B(U_M)}{A_0(U_M)}&=\frac{m}{m+q}d+o(d).
       \label{eq:r10-equipartition-law}
\end{align}
The distance includes translations and constant phases, already
contained in $\mathcal G_{\rm sel}$. For the multiplier in
\eqref{eq:r10-multiplier},
\begin{equation}\label{eq:r10-frequency-law}
 \omega_M=\frac{r}{M_0}\beta_*^m d^{m/q}(1+o(1)).
\end{equation}
\end{theorem}

\begin{proof}
For the normalization $M(V)=M_0$, $A_0(V)=1$, let
$H(V)=1-(2/p)\|V\|_p^p\ge0$. The exact energy formula is
\begin{equation}\label{eq:r10-exact-scaled-energy}
 I(M)=\frac{\alpha\lambda_M^m}{2}
          [-d+(1+d)H(V_M)]
       +\frac\alpha2\int b(\lambda_M\xi)|\widehat V_M(\xi)|^2\,d\xi.
\end{equation}
Negativity implies $0\le H(V_M)<d/(1+d)$. Therefore normalized
sharp-GN compactness applies, giving subsequential convergence, modulo
translation and phase, to some $Q\in\mathcal G_0$ strongly in $H^s$.

First $\lambda_M\to0$. Otherwise along a subsequence
$\lambda_M\ge\ell>0$, and \eqref{eq:r10-bounds} gives
\[
 \frac{B(U_M)}{A_0(U_M)}
 \ge c\int\min\{\lambda_M^q|\xi|^{m+q},|\xi|^m\}
                         |\widehat V_M(\xi)|^2\,d\xi.
\]
Choose a compact annulus carrying positive Fourier mass of the limiting
$Q$. The right side is then bounded below by a fixed positive number.
But \eqref{eq:r10-exact-scaled-energy} implies $B/A_0<d\to0$, a
contradiction.

Define the rescaled correction form
\[
 D_\lambda(V)=\int\frac{b(\lambda\xi)}{\lambda^{m+q}}
                              |\widehat V(\xi)|^2\,d\xi.
\]
Whenever $\lambda_n\downarrow0$ and $V_n\to Q$ strongly in $H^s$,
\begin{equation}\label{eq:r10-correction-liminf}
 \liminf_n D_{\lambda_n}(V_n)\ge C_*(Q).
\end{equation}
Indeed on a compact annulus, \eqref{eq:r10-low-expansion} gives uniform
multiplier convergence and strong Fourier $L^2$ convergence gives
convergence of the integral. Discard the complement by nonnegativity,
then exhaust all frequencies. In particular the liminf is at least
$C_{\min}$ for the normalized minimizing profiles. This step does not
assume uniform higher moments of those profiles.

For a fixed $Q\in\mathcal G_{\rm sel}$,
\eqref{eq:r10-bounds} bounds the rescaled multiplier by
$C|\xi|^{m+q}$. Dominated convergence gives
$D_\lambda(Q)\to C_{\min}$. Testing with
$\sqrt\alpha(\beta_*d^{1/q})^{N/2}Q(\beta_*d^{1/q}x)$ proves
\begin{equation}\label{eq:r10-energy-limsup}
 \limsup_{M\downarrow M_0}\frac{I(M)}{\alpha d^{1+m/q}}
 \le-K_*<0.
\end{equation}
On the other hand, negativity in
\eqref{eq:r10-exact-scaled-energy} gives
$\lambda_M^qD_{\lambda_M}(V_M)<d$. Equation
\eqref{eq:r10-correction-liminf} consequently yields
$\lambda_M^q\le2d/C_{\min}$ for small $d$. The lower bound
$I(M)\ge-\alpha d\lambda_M^m/2$, combined with
\eqref{eq:r10-energy-limsup}, yields
$\lambda_M^m\ge K_*d^{m/q}$ for small $d$. Thus
\begin{equation}\label{eq:r10-scale-comparable}
 0<c\le\beta_M:=\lambda_M/d^{1/q}\le C<\infty.
\end{equation}
The same inequalities bound $D_{\lambda_M}(V_M)$ and $H(V_M)/d$.

Extract a subsequence with $\beta_M\to\beta\in(0,\infty)$,
$V_M\to Q\in\mathcal G_0$ strongly in $H^s$ modulo symmetries, and
convergent bounded scalar quantities $D_{\lambda_M}(V_M)$ and
$H(V_M)/d$. Divide \eqref{eq:r10-exact-scaled-energy} by
$\alpha d^{1+m/q}$:
\begin{equation}\label{eq:r10-final-scaled-energy}
 \frac{I(M)}{\alpha d^{1+m/q}}
 =\frac12\left[-\beta_M^m
 +(1+d)\beta_M^m\frac{H(V_M)}d
 +\beta_M^{m+q}D_{\lambda_M}(V_M)\right].
\end{equation}
Its liminf is at least
\[
 \tfrac12[-\beta^m+C_*(Q)\beta^{m+q}]
 \ge\min_{t>0}\tfrac12[-t^m+C_{\min}t^{m+q}]=-K_*.
\]
Equality with \eqref{eq:r10-energy-limsup} forces
$\beta=\beta_*$, the unique scalar minimizer, and
$C_*(Q)=C_{\min}$. It also forces
$H(V_M)/d\to0$ and $D_{\lambda_M}(V_M)\to C_{\min}$.
Every extracted limit has these properties. This proves
\eqref{eq:r10-energy-law}--\eqref{eq:r10-deficit-law}, including the
distance statement. Finally,
\[
 \frac{B(U_M)}{A_0(U_M)}
 =\lambda_M^qD_{\lambda_M}(V_M)
 =d\beta_*^qC_{\min}+o(d)=\frac{m}{m+q}d+o(d).
\]
The multiplier identity \eqref{eq:r10-multiplier}, together with
$A_0(U_M)=\alpha\lambda_M^m$, $B/A_0=O(d)$, and $I/A_0=O(d)$, gives
$\omega_M/\lambda_M^m\to r/M_0$, proving
\eqref{eq:r10-frequency-law}.
\end{proof}

The proof also gives the normalized frequency correction
\[
 \frac{\omega_M}{\lambda_M^m}
 =\frac r{M_0}+\frac1{M_0}\left(r+\frac q{m+q}\right)d+o(d).
\]
This does not assert an additional-order expansion of $\lambda_M$.

\subsection{The limiting functional}

The selection has a precise variational interpretation. For $d>0$ set
$\alpha_d=(1+d)^{1/r}$ and $\ell=d^{1/q}$, and define, on $H^s$,
\[
 \mathcal F_d(W)=
 \frac{E_k(\sqrt{\alpha_d}\,\ell^{N/2}W(\ell\,\cdot))}
      {\alpha_d d^{1+m/q}}
 \quad\hbox{if }M(W)=M_0,
 \qquad \mathcal F_d(W)=+\infty\quad\hbox{otherwise}.
\]
The exact expression on the mass constraint is
\begin{equation}\label{eq:r10-gamma-exact}
 \mathcal F_d(W)=\frac{E_0(W)}d-\frac1p\|W\|_p^p
       +\frac12\int\frac{b(\ell\xi)}{\ell^{m+q}}
                                 |\widehat W(\xi)|^2\,d\xi.
\end{equation}

\begin{proposition}[Variational selection on the equality manifold]
\label{prop:r10-gamma}
In the strong $H^s$ topology, $\mathcal F_d$ $\Gamma$-converges as
$d\downarrow0$ to
\[
 \mathcal F_0(W)=
 \begin{cases}
 \tfrac12[C_*(W)-A_0(W)],&M(W)=M_0,\ E_0(W)=0,\\
 +\infty,&\text{otherwise}.
 \end{cases}
\]
Its minimizers are exactly the dilates
$\beta_*^{N/2}Q(\beta_*\,\cdot)$ with
$Q\in\mathcal G_{\rm sel}$, including their symmetries.
\end{proposition}

\begin{proof}
On the fixed mass constraint $E_0\ge0$. If $W_n\to W$ strongly in
$H^s$ and $E_0(W)>0$, continuity of the homogeneous energy and the
nonlinear norm makes \eqref{eq:r10-gamma-exact} tend to $+\infty$.
If $E_0(W)=0$, discard its nonnegative first term and apply the annular
liminf argument \eqref{eq:r10-correction-liminf}. Since equality in
the GN inequality gives $\|W\|_p^p/p=A_0(W)/2$, this is the required
liminf. A limit of the fixed mass constraint has mass $M_0$ by strong
$L^2$ convergence. For finite $\mathcal F_0(W)$, the constant sequence
$W$ is a recovery sequence, using domination by
$C|\xi|^{m+q}|\widehat W|^2$. This proves the $\Gamma$-limit.
Every nonzero equality profile is a dilation of an element of
$\mathcal G_0$; minimizing its scalar dilation gives exactly
\eqref{eq:r10-selection-constants}.
\end{proof}

Translation invariance prevents global equicoercivity in $H^s$; none
is asserted. Theorem~\ref{thm:r10-selection} separately proves
compactness of the rescaled minimizers modulo translations. The
limiting principle says that the first nonzero low-frequency correction
selects a subset of the entire homogeneous optimizer manifold.

\subsection{A joint small-perturbation and mass-excess limit}

There is a complementary selection limit that needs no higher-moment
hypothesis. Fix a positive elliptic homogeneous symbol $a$ as above and
let
\[
 b_0(\xi)=a(\xi)\frac{|\xi|^q}{1+|\xi|^q},\qquad
 k_\varepsilon=a+\varepsilon b_0,\qquad M_c=M_{c,a}.
\]
Write $A=A_a$, $B_0(f)=\int b_0|\widehat f|^2$, and let
$I_\varepsilon$ be the constrained infimum for $k_\varepsilon$. The
entire homogeneous equality manifold, with no kinetic normalization, is
\[
 \mathcal H=\{Q:M(Q)=M_c,\ E_a(Q)=0\}.
\]
For $0<\tau<1$ define
\[
 J(\tau)=\inf_{Q\in\mathcal H}\tfrac12[B_0(Q)-\tau A(Q)],\qquad
 \mathcal S_\tau=\operatorname*{argmin}_{Q\in\mathcal H}
                                \tfrac12[B_0(Q)-\tau A(Q)].
\]

\begin{proposition}[Selection by an arbitrarily weak perturbation]
\label{prop:r10-crossover}
For every $0<\tau<1$, $J(\tau)$ is finite, negative, and attained, and
$\mathcal S_\tau$ is compact in $H^s$ modulo translations and phase.
Let $\varepsilon\downarrow0$, set
\[
 \alpha_\varepsilon=(1+\varepsilon\tau)^{1/r},\qquad
 M_\varepsilon=\alpha_\varepsilon M_c,
\]
and let $U_\varepsilon$ be any minimizer of
$I_\varepsilon(M_\varepsilon)$. Then
\[
 \frac{I_\varepsilon(M_\varepsilon)}
             {\alpha_\varepsilon\varepsilon}\longrightarrow J(\tau),
 \qquad
 \operatorname{dist}_{H^s}
       (U_\varepsilon/\sqrt{\alpha_\varepsilon},\mathcal S_\tau)
       \longrightarrow0.
\]
No spatial rescaling is used. The same assertions hold along
$\tau_\varepsilon\to\tau\in(0,1)$, with
$\alpha_\varepsilon=(1+\varepsilon\tau_\varepsilon)^{1/r}$ and limiting
set $\mathcal S_\tau$; the compactness estimates are uniform for
$\tau_\varepsilon$ in a compact subinterval of $(0,1)$.
\end{proposition}

\begin{proof}
For every $\eta>0$,
$B_0(f)\ge(1-\eta)A(f)-C_\eta M(f)$. Thus the objective defining
$J(\tau)$ is coercive in $A$ on $\mathcal H$ if $\eta<1-\tau$.
Low-frequency dilation of any homogeneous optimizer makes the
objective negative. Its minimizing sequences consequently have $A$
bounded above and bounded away from zero, the latter because the
objective is at least $-\tau A/2$. Normalize $A$ to one and apply
sharp-GN compactness. The normalizing scales stay in a compact subset
of $(0,\infty)$, giving strong $H^s$ convergence modulo translations.
Continuity of the order-$m$ form $B_0$ proves attainment and compactness
of the selected set.

On the mass constraint $M(f)=M_c$, the exact identity is
\begin{equation}\label{eq:r10-crossover-functional}
 \frac{E_{k_\varepsilon}(\sqrt{\alpha_\varepsilon}f)}
             {\alpha_\varepsilon\varepsilon}
 =\frac{E_a(f)}\varepsilon+\frac12B_0(f)
                              -\frac\tau p\|f\|_p^p.
\end{equation}
It $\Gamma$-converges in the strong $H^s$ topology to
$[B_0-\tau A]/2$ on $\mathcal H$ and $+\infty$ off $\mathcal H$:
use $E_a\ge0$, continuity of the last two terms, and constant equality
profiles as recovery sequences.

We prove compactness of actual minimizers separately. The sharp GN
inequality at mass $\alpha_\varepsilon M_c$ gives
\[
 E_{k_\varepsilon}(U)
 \ge\frac\varepsilon2
       [(1-\eta-\tau)A(U)-C_\eta M(U)].
\]
Testing \eqref{eq:r10-crossover-functional} with a selected profile
gives the exact bound
$I_\varepsilon/(\alpha_\varepsilon\varepsilon)\le J(\tau)<0$.
Hence $A(U_\varepsilon)$ is uniformly bounded. The other lower bound
$E_{k_\varepsilon}(U)\ge-\varepsilon\tau A(U)/2$ forces it to be
bounded away from zero. With
$f_\varepsilon=U_\varepsilon/\sqrt{\alpha_\varepsilon}$, identity
\eqref{eq:r10-crossover-functional} now gives
$E_a(f_\varepsilon)=O(\varepsilon)$. Normalize the kinetic form and
use sharp-GN compactness to obtain a strong $H^s$ limit
$Q\in\mathcal H$ modulo translation, without any escaping spatial
scale. The $\Gamma$-liminf is at least $J(\tau)$ and the trial upper
bound is at most $J(\tau)$, proving convergence of the energies and
selection of $Q\in\mathcal S_\tau$.

For $\tau_\varepsilon\to\tau$, the same proof has uniform constants:
choose $\eta$ below the distance of the compact parameter interval to
one, and a single sufficiently low-frequency optimizer trial that is
strictly negative for the smallest parameter in that interval.
Continuity of the bounded-energy objective gives the stated limit.
\end{proof}

Along every convergent subsequence in this proposition, the multiplier
satisfies $\omega_\varepsilon\to rA(Q)/M_c$ for
$Q\in\mathcal S_\tau$. Thus a perturbation tending to zero can select
a finite spatial scale and frequency from the critical homogeneous
family when the mass excess has the same order as the perturbation.
The selection functional $J$ is concave and strictly decreasing. The
first property follows by taking an infimum of affine functions; for
the second, test the objective at $\tau_2>\tau_1$ with a minimizer
$Q_1$ at $\tau_1$ to obtain
\[
 J(\tau_2)\le J(\tau_1)-\tfrac12(\tau_2-\tau_1)A(Q_1)<J(\tau_1).
\]

\subsection{A quartic example with explicit exponents}

For $a_0(\xi)=|\xi|^4$ and
$b(\xi)=\varepsilon|\xi|^6/(1+|\xi|^2)$, let
\[
 S_6=\min_{Q\in\mathcal G_0}\int|\xi|^6|\widehat Q(\xi)|^2\,d\xi,
 \qquad d=(M/M_0)^{4/N}-1,\qquad \alpha=M/M_0.
\]
The higher-moment hypothesis holds, and $C_{\min}=\varepsilon S_6$.
Theorem~\ref{thm:r10-selection} gives
\[
 I(M)\sim-\alpha\frac{2}{27\varepsilon^2 S_6^2}d^3,
 \qquad \lambda_M^2\sim\frac{2}{3\varepsilon S_6}d,
 \qquad
 \omega_M\sim\frac{4}{NM_0}\frac{4}{9\varepsilon^2 S_6^2}d^2.
\]
Thus the spatial dilation length grows as $d^{-1/2}$, the principal
kinetic energy vanishes as $d^2$, and the negative constrained energy
vanishes as $d^3$. The selected critical profiles minimize their sixth
Fourier moment. At the other end, $q=2<m=4$ implies
$I(M_\infty)=-\infty$.

For the isotropic family as a whole the result is a complete constrained
energy classification: zero nonattained infimum up to $M_0$, negative
attained infimum in $(M_0,M_\infty)$, and $-\infty$ above
$M_\infty$, with the upper endpoint determined by the three cases in
Theorem~\ref{thm:r10-endpoint}. The new asymptotic concerns a class of
nonhomogeneous equations; it does not resolve the unrestricted dynamics
of the original homogeneous BNLS.

\section{Ballistic exterior-gradient escape and weaker localization rigidity}
\label{sec:r10-dynamics}

We consider the pure mass-critical equation
\begin{equation}\label{eq:r10dyn-equation}
 iu_t=\Delta^2u-|u|^{8/N}u,\qquad N\ge1,
\end{equation}
with the conventions
\[
 p=2+8/N,\quad M=\|u\|_2^2,\quad
 P=\operatorname{Im}\int\overline u\nabla u,\quad
 E=\tfrac12\|\Delta u\|_2^2-\tfrac1p\|u\|_p^p.
\]
Throughout this section, a solution is a standard strong $H^2$ solution
for which these quantities are conserved and the localized virial identity
holds in its integrated energy-space form. Suppose it is forward global and
\begin{equation}\label{eq:r10dyn-bound}
 \sup_{t\ge0}\|u(t)\|_{H^2}\le B<\infty.
\end{equation}
No radial symmetry or finite spatial moment is imposed. Our main result is a
necessary spatial-spreading condition for a hypothetical bounded global
negative-energy solution with zero momentum. It does not establish that
such solutions exist, and it does not exclude all such solutions.

\subsection{A tail estimate and a moving-center identity}

For a center $c(t)$, write
\[
 \begin{gathered}
 D_R^c(t)=\int_{|x-c(t)|>R}|\nabla u(t,x)|^2\,dx,\\
 q_N=\min(4,16/N),\quad
 \kappa_N=\min(1/2,8/N),\quad \beta_N=\min(2,q_N).
 \end{gathered}
\]
The distinction between gradient localization and mass localization is
useful: the former allows noncompact $L^2$ tails whose gradients vanish.

\begin{lemma}[Gradient control of the nonlinear tail]
\label{lem:r10dyn-nonlinear}
For every center $c(t)$ and $R\ge1$,
\begin{equation}\label{eq:r10dyn-nonlineartail}
 \int_{|x-c(t)|>2R}|u(t,x)|^p\,dx
 \le C_{N,M,B}\bigl(D_R^c(t)^{1/2}+R^{-1}\bigr)^{q_N}.
\end{equation}
If $E<0$, there is also $g_*>0$, depending only on $N,M,B,|E|$,
such that $\|\nabla u(t)\|_2\ge g_*$ for all $t\ge0$.
\end{lemma}
\begin{proof}
For $f\in H^2$, homogeneous Sobolev embedding and interpolation give
\begin{equation}\label{eq:r10dyn-interpolation}
 \|f\|_p^p\le
 \begin{cases}
 C_N\|f\|_2^{2(4-N)/N}\|\nabla f\|_2^4,&N\le4,\\[2pt]
 C_N\|\nabla f\|_2^{16/N}\|\Delta f\|_2^{2-8/N},&N>4.
 \end{cases}
\end{equation}
Indeed, the derivative order in the homogeneous embedding into $L^p$ is
$\sigma=N(1/2-1/p)=2N/(N+4)$. Interpolate between orders zero and one
when $\sigma\le1$, and between orders one and two when $\sigma>1$.
This includes the usual $\dot H^1\hookrightarrow L^4$ endpoint at $N=4$.

Choose a smooth cutoff $\eta$ equal to zero on the unit ball and to one
outside the ball of radius two. For
$f(x)=\eta((x-c(t))/R)u(t,x)$, product differentiation gives
\[
 \|f\|_2\le\sqrt M,\quad
 \|\nabla f\|_2\le D_R^c(t)^{1/2}+C\sqrt M/R,\quad
 \|\Delta f\|_2\le C_B.
\]
Substitution in \eqref{eq:r10dyn-interpolation} proves
\eqref{eq:r10dyn-nonlineartail}. Finally, negative energy implies
$\|u(t)\|_p^p=p\|\Delta u(t)\|_2^2/2+p|E|\ge p|E|$.
The global inequalities \eqref{eq:r10dyn-interpolation}, mass conservation,
and \eqref{eq:r10dyn-bound} therefore give the positive gradient lower bound.
\end{proof}

Choose a smooth radial function $\phi_R$ which equals $|x|^2/2$ for
$|x|\le4R$, is constant for $|x|\ge8R$, and satisfies
\[
 D^2\phi_R\le I,\qquad
 \|D^j\phi_R\|_\infty\le C_jR^{2-j}\quad(1\le j\le6).
\]
For a Lipschitz center $a$ with $\|a'\|_\infty\le L$, define
\begin{equation}\label{eq:r10dyn-virialfunctional}
 J_R(t)=2\operatorname{Im}\int
 \nabla\phi_R(x-a(t))\cdot\overline u(t,x)\nabla u(t,x)\,dx.
\end{equation}

\begin{lemma}[Moving virial with gradient-tail error]
\label{lem:r10dyn-virial}
There are constants $C_0=C_0(N,M,B)>0$ and
$C_1=C_1(N,M,B,L)>0$ such that
\begin{align}
 |J_R(t)|&\le C_0R,\label{eq:r10dyn-endpoint}\\
 J_R'(t)&\le16E-2a'(t)\cdot P+
 C_1\bigl(R^{-\beta_N}+D_R^a(t)^{\kappa_N}\bigr).
 \label{eq:r10dyn-virialbound}
\end{align}
The derivative inequality is understood in the integrated sense.
\end{lemma}
\begin{proof}
The endpoint estimate follows directly from Cauchy--Schwarz. Put
$j=\operatorname{Im}(\overline u\nabla u)$ and $\alpha=8/N$.
At a fixed center, differentiation and integration by parts give
\begin{align*}
 J_R'={}&8\operatorname{Re}\int\partial_{jk}\phi_R\,
      \partial_{ij}\overline u\,\partial_{ik}u
 -4\operatorname{Re}\int\partial_{jk}\Delta\phi_R\,
      \partial_j\overline u\,\partial_ku\\
 &-2\int\Delta^2\phi_R|\nabla u|^2
 +\int\Delta^3\phi_R|u|^2
 -\frac{2\alpha}{p}\int\Delta\phi_R|u|^p.
\end{align*}
Repeated indices are summed, and the cutoff is evaluated at $x-a(t)$.
Since $D^2\phi_R\le I$, the leading term is at most
$8\|D^2u\|_2^2=8\|\Delta u\|_2^2$. The higher cutoff derivatives
contribute $O_B(R^{-2})$. The nonlinear term equals
$-16\|u\|_p^p/p$ up to an upper error bounded by
$C\int_{|x-a(t)|>4R}|u|^p$. These principal terms sum to $16E$.

Motion of the center adds exactly
\[
 -2a'(t)\cdot\int D^2\phi_R(x-a(t))j(t,x)\,dx.
\]
Replacing the Hessian by $I$ gives $-2a'\cdot P$. Its error is at most
\[
 C L\int_{|x-a(t)|>4R}|u||\nabla u|
 \le C L\sqrt M\,D_{4R}^a(t)^{1/2}.
\]
Apply Lemma~\ref{lem:r10dyn-nonlinear}, using
$D_{2R}^a,D_{4R}^a\le D_R^a$ and $0\le D_R^a\le B^2$.
The nonlinear-tail exponent is $q_N/2=\min(2,8/N)$; the
center error has exponent $1/2$. Reducing both to
$\kappa_N$ on this bounded interval proves
\eqref{eq:r10dyn-virialbound}. The calculation passes from smooth solutions
to the integrated $H^2$ identity using the displayed bounds. A Lipschitz
center has an almost-everywhere derivative; alternatively, smooth the
center and pass to the integrated identity.
\end{proof}

\subsection{A ballistic necessary alternative}

\begin{theorem}[Persistent exterior gradient at a linear distance]
\label{thm:r10dyn-ballistic}
Suppose \eqref{eq:r10dyn-bound} holds, $E=-e<0$, and $P=0$.
For every Lipschitz center $a:[0,\infty)\to\mathbb R^N$ with
$\|a'\|_\infty\le L$, there are constants $\nu,d>0$, depending only
on $N,M,B,e,L$, such that
\begin{equation}\label{eq:r10dyn-ballistic}
 \liminf_{T\to\infty}\inf_{S\ge0}\frac1T
 \int_S^{S+T}\int_{|x-a(t)|>\nu T}|\nabla u(t,x)|^2\,dx\,dt
 \ge d.
\end{equation}
In particular no choice of bounded-speed center can localize the gradient
inside a sublinearly growing spatial radius in this averaged sense.
\end{theorem}
\begin{proof}
Integrate \eqref{eq:r10dyn-virialbound} on $[S,S+T]$, use $P=0$,
and apply \eqref{eq:r10dyn-endpoint}. Uniformly in $S$,
\begin{equation}\label{eq:r10dyn-integratedescape}
 C_1\frac1T\int_S^{S+T}D_R^a(t)^{\kappa_N}\,dt
 \ge16e-C_1R^{-\beta_N}-2C_0R/T.
\end{equation}
Choose a fixed $\nu>0$ so small that $2C_0\nu\le4e$, and put
$R=\nu T$, held constant during each time integration.
For all sufficiently large $T$, the right side is at least
$8e$. Since $0<\kappa_N\le1$, Jensen's inequality gives
\[
 \left(\frac1T\int_S^{S+T}D_{\nu T}^a(t)\,dt\right)^{\kappa_N}
 \ge\frac1T\int_S^{S+T}D_{\nu T}^a(t)^{\kappa_N}\,dt
 \ge8e/C_1.
\]
This proves \eqref{eq:r10dyn-ballistic} with
$d=(8e/C_1)^{1/\kappa_N}$. The argument also proves the assertion
uniformly over all centers with the specified speed bound.
\end{proof}

\begin{remark}
The theorem asserts persistent exterior gradient in a time-averaged sense,
not at every time. Since $D_R^a(t)\le B^2$, the lower bound also forces a
positive fraction of times in every sufficiently long window to have a
nonzero exterior-gradient amount. The radius $\nu T$ is proportional to
the length of that window. No uniform positive $\nu$ or $d$ is claimed
as $E\to0$ or $B\to\infty$. The exterior gradient is not a physical
dissipation, and the theorem supplies no finite budget for its repeated
transport or replenishment.
\end{remark}

\subsection{Uniform and averaged gradient-localization rigidity}

\begin{theorem}[Averaged localization forces counter-momentum drift]
\label{thm:r10dyn-cesaro}
Suppose \eqref{eq:r10dyn-bound} holds and $E<0$. Let $a$ be Lipschitz,
and suppose
\begin{equation}\label{eq:r10dyn-cesarohyp}
 \lim_{R\to\infty}\limsup_{T\to\infty}\frac1T
 \int_0^TD_R^a(t)\,dt=0.
\end{equation}
Then
\begin{equation}\label{eq:r10dyn-drift}
 \limsup_{T\to\infty}\frac{P\cdot[a(T)-a(0)]}{T}\le8E.
\end{equation}
Consequently \eqref{eq:r10dyn-cesarohyp} is impossible when $P=0$.
\end{theorem}
\begin{proof}
Integrating \eqref{eq:r10dyn-virialbound} and rearranging gives
\[
 2\frac{P\cdot[a(T)-a(0)]}{T}
 \le16E+C_1R^{-\beta_N}+
 C_1\left(\frac1T\int_0^TD_R^a(t)\,dt\right)^{\kappa_N}
 +2C_0R/T.
\]
Here Jensen was again used. First take the upper limit as $T\to\infty$
for fixed $R$, and then let $R\to\infty$. This proves the result.
\end{proof}

\begin{lemma}[Regularizing a gradient concentration center]
\label{lem:r10dyn-center}
Suppose \eqref{eq:r10dyn-bound} holds, $E<0$, and an arbitrary map
$c:[0,\infty)\to\mathbb R^N$ satisfies
\begin{equation}\label{eq:r10dyn-uniformtight}
 \sup_{t\ge0}D_R^c(t)\longrightarrow0\qquad(R\to\infty).
\end{equation}
Then there is a Lipschitz map $a$ with
$\sup_{t\ge0}|a(t)-c(t)|<\infty$. Gradient localization also holds
about $a$.
\end{lemma}
\begin{proof}
The nonlinearity maps bounded $H^2$ sets into bounded $H^{-2}$ sets:
for a test function $h\in H^2$, H\"older and $H^2\hookrightarrow L^p$ give
\[
 \left|\int |u|^{p-2}u\overline h\right|
 \le\|u\|_p^{p-1}\|h\|_p\le C_B\|h\|_{H^2}.
\]
The equation therefore implies a uniform bound on $u_t$ in $H^{-2}$.
Interpolation between $H^{-2}$ and the bounded $H^2$ norm yields
\begin{equation}\label{eq:r10dyn-timeholder}
 \|u(t)-u(s)\|_{H^1}\le C_B|t-s|^{1/4}.
\end{equation}
Let $g_*>0$ be the gradient lower bound from
Lemma~\ref{lem:r10dyn-nonlinear}. Choose $R_0$ so that
$\sup_tD_{R_0}^c(t)\le g_*^2/16$, and choose $\tau>0$ so that
\eqref{eq:r10dyn-timeholder} is at most $g_*/4$ when $|t-s|\le\tau$.

If $|c(t)-c(s)|>2R_0$, the two radius-$R_0$ balls are disjoint.
On $B(c(s),R_0)$, the gradient of $u(s)$ has $L^2$ norm at least
$\sqrt{15}g_*/4$, whereas that of $u(t)$ has norm at most $g_*/4$.
The reverse triangle inequality contradicts
\eqref{eq:r10dyn-timeholder}. Thus
$|c(t)-c(s)|\le2R_0$ whenever $|t-s|\le\tau$.

Let $a$ interpolate linearly through the values $c(n\tau)$,
$n=0,1,\ldots$. Then
\[
 \|a'\|_\infty\le2R_0/\tau,\qquad
 \sup_{t\ge0}|a(t)-c(t)|\le4R_0.
\]
The bounded displacement preserves the uniform tail condition.
\end{proof}

\begin{corollary}[Uniform gradient tightness, with no center regularity]
\label{cor:r10dyn-uniform}
Under \eqref{eq:r10dyn-bound}, $E<0$, and
\eqref{eq:r10dyn-uniformtight},
\[
 \limsup_{T\to\infty}\frac{P\cdot[c(T)-c(0)]}{T}\le8E.
\]
In particular there is no forward-global uniformly $H^2$-bounded
negative-energy zero-momentum solution with uniformly localized gradient
modulo translations. This holds at every mass.
\end{corollary}
\begin{proof}
Apply Lemma~\ref{lem:r10dyn-center}, followed by
Theorem~\ref{thm:r10dyn-cesaro}. The bounded difference between the two
centers does not change the limiting drift quotient.
\end{proof}

Uniform mass tightness and \eqref{eq:r10dyn-bound} imply uniform gradient
tightness by cutoff interpolation. The converse is false for general
bounded $H^2$ families: a fixed localized core can be accompanied by a
fixed-mass packet whose width tends to infinity and whose gradient norm
tends to zero. The result above consequently weakens the required spatial
localization, although it does not weaken the uniform-in-time $H^2$ bound.

\subsection{The unresolved near-critical step}

These results identify an actual requirement on any prospective bounded
global counterexample, but they do not remove the general dynamical gap.
If $M_c$ is the sharp GN critical mass and
$\delta_M=(M/M_c)^{4/N}-1$, the sharp inequality gives
\[
 -\tfrac12\delta_M A(u)\le E(u)<0,\qquad
 A(u)=\|\Delta u\|_2^2.
\]
For fixed negative energy, $A\ge2|E|/\delta_M$. Conversely, normalizing
$A=1$ along a sequence $M\downarrow M_c$ forces $E\to0$, so compactness
can produce an allowed zero-energy optimizer. Normalizing $E=-1$ instead
forces the kinetic norm to diverge. Near-optimizer compactness therefore
does not, by itself, preserve a negative-energy contradiction at the
limiting mass.

A near-critical orbit may have a small, nonzero exterior-gradient budget;
$|E|/A$ can be arbitrarily smaller than that budget. The positive lower
bounds in Theorem~\ref{thm:r10dyn-ballistic} can be compatible with this
small tail. Closing the unrestricted problem requires a further mechanism
that controls sustained radiation or its replenishment, or a construction
showing that such behavior is possible. No such additional mechanism is
asserted here.


\section{Relation to established methods and remaining questions}
Concentration compactness, translation profile decomposition, sharp
interpolation, constrained minimization and localized virial identities
are established tools. The claims requiring mathematical and historical
assessment here are the particular threshold compactification,
spectral bounds, uniform nonlocal gap, and dispersion-dependent
selection statements in their precise stated classes.

Pausader--Shao~\cite{PS} already use compact-orbit virial and center-motion
arguments for the mass-critical fourth-order equation. Our gradient-tail
and ballistic conclusions are consequences of that virial framework
with the displayed localization estimates; we do not present the virial
identity itself as new. The dynamical result allows noncompact mass
tails, but it still assumes a uniform energy-space bound on the orbit.

Bonheure--Cast\'eras--Gou--Jeanjean~\cite{BCGJ} analyze normalized
solutions for mixed second- and fourth-order dispersion in a critical
and supercritical regime. Liu--Romani--Su~\cite{LRS} study relations
between optimizers and ground states in the mass-competition regime
of that mixed-order equation. Those works supply relevant variational
context. The same-order low/high frequency interpolation, its endpoint
trichotomy and its optimizer-manifold selection laws are the specific
objects analyzed here; a targeted comparison does not prove priority.

Cho--Hwang--Kwon--Lee~\cite{CHKL}, Theorems 1.1 and 1.2, prove
negative-energy finite-time blowup for fractional/higher-order Hartree
equations under radial, high-regularity and spatially weighted initial
data assumptions. Their higher-order theorem has
$2<\alpha<1+N/2$. It is therefore not an unrestricted nonradial theorem
for every symbol and interaction in Section~\ref{sec:r10-hartree}.
The Hartree result here is a traveling-profile exclusion theorem; no
contradiction with that blowup theorem is asserted.

The central unfinished dynamical question in this continuation remains
the existence or exclusion of a forward-global uniformly $H^2$-bounded
pure BNLS orbit with $P=0$, $E<0$, and mass in a fixed sufficiently
small interval above $M_c$. The new escape theorem identifies a necessary
mechanism for such an orbit but does not make that mechanism impossible.
Even excluding all global bounded orbits would not, by itself, exclude
infinite-time growth of the $H^2$ norm or prove finite-time blowup.

For the static threshold, the next precise questions are whether the
endpoint has negative or zero energy, whether more than one component
can occur, and which spectral degeneracies permit the endpoint to be
approached. For modified dispersion the full isotropic constrained
classification is proved, while general anisotropic borderline
upper-endpoint equality and nonlinear stability of individual branches
require additional analysis. These are separate mathematical questions;
none is settled merely by increasing the number of parameters covered.

\begin{thebibliography}{9}
\bibitem{PS} B. Pausader and S. Shao,
\emph{The mass-critical fourth-order Schr\"odinger equation in high dimensions},
J. Hyperbolic Differ. Equ. \textbf{7} (2010), 651--705,
\href{https://arxiv.org/abs/0906.1547}{arXiv:0906.1547}.
\bibitem{BCGJ} D. Bonheure, J.-B. Cast\'eras, T. Gou and L. Jeanjean,
\emph{Normalized solutions to the mixed dispersion nonlinear Schr\"odinger
equation in the mass critical and supercritical regime},
\href{https://arxiv.org/abs/1802.09217}{arXiv:1802.09217v2}.
\bibitem{LRS} Z. Liu, G. Romani and Y. Su,
\emph{Mixed-dispersion Schr\"odinger equations and Gagliardo--Nirenberg
inequalities: equivalence between ground states and optimizers},
\href{https://arxiv.org/abs/2603.26493}{arXiv:2603.26493v1} (2026).
\bibitem{CHKL} Y. Cho, G. Hwang, S. Kwon and S. Lee,
\emph{On finite time blowup for the mass-critical Hartree equations},
\href{https://arxiv.org/abs/1208.2302}{arXiv:1208.2302v3} (2014).
\bibitem{LZ} Z. Lin and C. Zeng,
\emph{Instability, index theorem, and exponential trichotomy for linear
Hamiltonian PDEs}, Mem. Amer. Math. Soc. \textbf{275} (2022), no. 1347,
\href{https://arxiv.org/abs/1703.04016}{arXiv:1703.04016v4}.
\end{thebibliography}
\end{document}
